\textsc{We assume the reader's familiarity}
with basic category theoretical concepts,
as well as the main features of the theories of monoidal and 2-categories,
as discussed for example in~\cite{MacLane1998,Etingof2015,category-theory-in-context_riehl,johnson-yau2021:TwoDimensionalCategories}.
Nevertheless, we start with a brief review of the most important terminology and notation.

\section{Bicategories}\label{sec:bicategories}

\textsc{The theory of bicategories is essential}
to the graphical calculus we will use throughout,
see \cref{sec:graphical-calculus-introduction},
as well as to some formal arguments we shall employ.
Following~\cite{johnson-yau2021:TwoDimensionalCategories},
we recall the most important definitions.

\begin{definition}\label{def:bicategory}
  A \emph{bicategory} \(\mathbb{B}\) has as data:
  \begin{itemize}
    \item A collection of \emph{objects} \(\Ob\mathbb{B}\), where we often write \(x \in \mathbb{B}\) for \(x \in \Ob \mathbb{B}\);
    \item for all \(x, y \in \mathbb{B}\) a \emph{hom-category} \(\mathbb{B}(x,y)\),
    whose objects are called \emph{1-cells} (or 1-mor\-phisms) and whose morphisms are called \emph{2-cells} (or 2-morphisms);
    \item for all \(x, y, z \in \mathbb{B} \) a \emph{horizontal composition}
    functor
    \[
      \otimes_y \from \mathbb{B}(y, z) \times \mathbb{B}(x, y) \to \mathbb{B}(x, z);
    \]
    \item for all \(x \in \mathbb{B}\) an \emph{identity} \(1_x \in \mathbb{B}(x,x) \);
    \item a natural isomorphism \(\alpha_{M, N, P} \from (M \otimes_y N) \otimes_x P \to M \otimes_y (N \otimes_x P)\);
    \item natural isomorphisms \(\lambda_M \from 1_y \otimes_y M \to M\) and \(\rho_M \from M \to M \otimes_x 1_x\).
  \end{itemize}

  \noindent{}This data is subject to the commutativity of the following diagrams, for all suitable 1-cells \(M, N, P, R\):
  \[
    \mathscale{0.9}{\begin{mytikzcd}[ampersand replacement=\&]
      {(M \otimes_z (N \otimes_y P)) \otimes_x R} \&\& {M \otimes_z ((N \otimes_y P) \otimes_x R)} \\
      {((M \otimes_z N) \otimes_y P) \otimes_x R} \&\& {M \otimes_z (N \otimes_y (P \otimes_x R))} \\
      \& {(M \otimes_z N) \otimes_y (P \otimes_x R)} \\
      {(M \otimes_x 1_x) \otimes_x N} \&\& {M \otimes_x (1_x \otimes_x N)} \\
      {M \otimes_x N} \&\& {M\otimes_x N}
      \arrow["{{\alpha_{M, N \otimes_y P, R}}}", from=1-1, to=1-3]
      \arrow["{{M \otimes \alpha_{NPR}}}", from=1-3, to=2-3]
      \arrow["{{\alpha_{MNP} \otimes R}}", from=2-1, to=1-1]
      \arrow["{{\alpha_{M \otimes_z N, P, R}}}"', from=2-1, to=3-2]
      \arrow["{{\alpha_{M,N,P \otimes_x R}}}"', from=3-2, to=2-3]
      \arrow["{{\alpha_{M, 1_x, N}}}", from=4-1, to=4-3]
      \arrow["{{M \otimes \lambda_N}}", from=4-3, to=5-3]
      \arrow["{{\rho_M \otimes N}}", from=5-1, to=4-1]
      \arrow[Rightarrow, no head, from=5-1, to=5-3]
    \end{mytikzcd}}
  \]

  If \(\alpha\), \(\lambda\), and \(\rho\) are all identities, we call \(\mathbb{B}\) a \emph{2-category}.
\end{definition}

The composition of morphisms inside of the category \(\mathbb{B}(x,y)\) will be called the \emph{vertical composition}.
We denote the terminal bicategory
consisting of a single object, a single 1-, and a single 2-morphism
by \(\terminal\).

\begin{example}\label{ex:cat-is-a-bicategory}
  The category \(\BiCat{Cat}\) of (small) categories, functors, and natural transformations is a 2-category.
  Horizontal composition is given by horizontal composition of functors and natural transformations.
\end{example}

Examples of bicategories are plentiful throughout this work.
We refrain from spelling them out here
before having introduced more of the necessary notation,
but see
\cref{%
  ex:bicategory-of-bimodules,%
  MonoidalBicatCat,%
  MonoidalBicatProf,%
  rmk:twisted-centre-as-loop,%
  rmk:bimonad-in-bicategorical-language%
}

\begin{remark}
  The use of \(\otimes\) for the horizontal composition is non-standard,
  and inspired by~\cite{garner22}.
  We use it here to highlight the parallels to the definition of monoidal categories,
  see \cref{def:monoidal-category}.
\end{remark}

\begin{notation}
  To syntactically differentiate bicategories from ordinary categories,
  we will generally start the name of a bicategory with a blackboard bold letter.
  For example, when talking about the 1-category of all (small) categories and functors,
  we write \(\Cat\) instead of \(\BiCat{Cat}\).
\end{notation}

Much like monoidal categories,
see \cref{thm:monoidal-strictness-theorem},
bicategories admit a \emph{coherence} and \emph{strictification} result,
in that every bicategory is biequivalent to a 2-category,
see for example~\cite[Theorem~8.4.1 and~Corollary~8.4.2]{johnson-yau2021:TwoDimensionalCategories}.

\begin{definition}\label{def:bifunctor}
  A \emph{lax functor} \(F\from \mathbb{B} \to \mathbb{C}\)
  of bicategories \(\mathbb{B}\), \(\mathbb{C}\)
  has as data:
  \begin{itemize}
    \item an object assignment \(F \from \Ob\mathbb{B} \to \Ob\mathbb{C}\);
    \item for all \(x,y \in\mathbb{B}\), a functor \(F_{x, y} \from \mathbb{B}(x, y) \to \mathbb{C}(F x, F y)\);
    \item a family of natural transformations \(F_2 \from F \otimes F \nt F \circ \otimes\),
    with components \(F_{2, M, N} \from F_{yz} M \otimes_{F y} F_{xy} N \to F_{xz} (M \otimes_y N)\),
    for all \(x,y,z \in\mathbb{B}\), \(M\in\mathbb{B}(y,z)\), and \(N \in\mathbb{B}(x,y)\); and
    \item for all \(x \in\mathbb{B}\), arrows \(F_{0, x} \from 1_{F x} \to F 1_x\).
  \end{itemize}

  \noindent{}This data is subject to the following axioms in \(\mathbb{C}(Fw, Fz)\) and \(\mathbb{C}(Fx, Fy)\),
  respectively, for all admissible \(M,N,P\):
  \[
    \begin{mytikzcd}[ampersand replacement=\&]
      {(FM \otimes_{Fy} FN) \otimes_{Fx} FP} \& {FM \otimes_{Fy} (FN \otimes_{Fx} FP)} \\
      {F(M \otimes_y N) \otimes_{Fx} FP} \& {FM \otimes_{Fy} F(N \otimes_x P)} \\
      {F((M \otimes_y N) \otimes_x P)} \& {F(M \otimes_y (N \otimes_x P))}
      \arrow["\alpha", from=1-1, to=1-2]
      \arrow["{F_{2; M, N} \otimes \mathrm{id}}"', from=1-1, to=2-1]
      \arrow["{\mathrm{id}\otimes F_{2; N, P}}", from=1-2, to=2-2]
      \arrow["{F_{2; M\otimes N, P}}"', from=2-1, to=3-1]
      \arrow["{F_{2; M, N \otimes P}}", from=2-2, to=3-2]
      \arrow["{F\alpha}"', from=3-1, to=3-2]
    \end{mytikzcd}
  \]
  \[
    \begin{mytikzcd}[ampersand replacement=\&]
      {1_{Fy} \otimes_{Fy} FM} \& FM \& {FM \otimes_{Fx} 1_{Fx}} \\
      {F1_y \otimes_{Fy} FM} \&\& {FM \otimes_{Fx} F1_x} \\
      {F(1_y \otimes _y M)} \& FM \& {F(M \otimes_x 1_x)}
      \arrow["\lambda", from=1-1, to=1-2]
      \arrow["{F_{0;y}\otimes\mathrm{id}}"', from=1-1, to=2-1]
      \arrow["\rho", from=1-2, to=1-3]
      \arrow[Rightarrow, no head, from=1-2, to=3-2]
      \arrow["{\mathrm{id}\otimes F_{0;x}}", from=1-3, to=2-3]
      \arrow["{F_{2;1_y,M}}"', from=2-1, to=3-1]
      \arrow["{F_{2;M,1_x}}", from=2-3, to=3-3]
      \arrow["{F\lambda}"', from=3-1, to=3-2]
      \arrow["{F\rho}"', from=3-2, to=3-3]
    \end{mytikzcd}
  \]

  A lax functor is called a \emph{pseudofunctor} if all \(F_2\)'s and \(F_0\)'s are invertible,
  and a \emph{2-functor} if they are all identities.
\end{definition}

Analogously, one defines \emph{oplax functors} between 2-categories.

\begin{definition}\label{def:oplax-transformation}
  An \emph{oplax transformation} \(\omega\from F \nt G\) between lax functors
  \((F, F_2, F_0), (G,G_2,G_0)\from \mathbb{B} \to \mathbb{C} \)
  consists of
  a 1-cell \(\omega_x \from Fx \to Gx\) in \(\mathbb{C}\) for every \(x \in \mathbb{B}\),
  and for all \(x, y \in \mathbb{B}\) a natural transformation
  \[
    \omega\from \omega_y \otimes F(\blank) \nt G(\blank) \otimes \omega_x\from \mathbb{B}(x,y) \to \mathbb{C}(Fx,Gy)
  \]
  with component 2-cells \(\omega_f\from \omega_y \otimes_{Fy} Ff \nt Gf \otimes_{Gx} \omega_x\) for all \(f\from x\to y\) in \(\mathbb{B}\).
  Graphically, we draw these 2-cells like the following:
  \[
    \begin{mytikzcd}[ampersand replacement=\&]
      Fx \& Fy \\
      Gx \& Gy
      \arrow["Ff", from=1-1, to=1-2]
      \arrow["{\omega_x}"', from=1-1, to=2-1]
      \arrow["{\omega_f}"', shorten <=8pt, shorten >=8pt, Rightarrow, from=1-2, to=2-1]
      \arrow["{\omega_y}", from=1-2, to=2-2]
      \arrow["Gf"', from=2-1, to=2-2]
    \end{mytikzcd}
  \]

  An oplax transformation is called a \emph{pseudonatural transformation}\note{%
    Pseudonatural transformations are called
    \emph{strong transformations} in~\cite{johnson-yau2021:TwoDimensionalCategories}.
  }
  if all \(\omega_f\)'s are isomorphisms,
  and a \emph{2-natural transformation} if they are all identities.
\end{definition}

Analogously, one defines \emph{lax} transformations between lax functors,
as well as lax and oplax transformations between oplax functors.

\begin{definition}[{\cite{lack10:icons}}]\label{def:icon}
  Let \(F, G\from\mathbb{B} \to\mathbb{C}\) be lax functors between bicategories.
  An \emph{\textsc{icon}} between \(F\) and \(G\) consists of
  an assertion that \(Fx = Gx\), for all \(x \in\mathbb{B}\); and
  an oplax transformation \(\omega \from F \nt G\),
  such that for all \(x \in \mathbb{B}\), the 1-cell \(\omega_x\from Fx \to Gx\) is the identity.
\end{definition}

\Cref{def:icon} might seem a bit contrived at first,
but it is an important concept:
there is no 2-category with
(small) bicategories as 0-cells,
lax (even pseudo-) functors as 1-cells,
and 2-natutral transformations as 2-cells.
Instead, \textsc{icon}s yield the ``right'' kind of 2-cells for this construction,
see~\cite[Section~4.6]{johnson-yau2021:TwoDimensionalCategories}
for a more extensive account.

\begin{definition}\label{def:modification}
  Let \( \mathbb{B}\) and \( \mathbb{C}\) be bicategories,
  \(F, G \from \mathbb{B} \to \mathbb{C}\) lax functors,
  and suppose that \(\alpha, \beta\from F \nt G\) are oplax transformations.
  A \emph{modification} \(\Upsilon \from \alpha \Rrightarrow \beta\) between \(\alpha\) and \(\beta\) consists of
  a 2-cell \(\Upsilon_x \from \alpha_x \nt \beta_x\) for every \(x \in \mathbb{B}\),
  such that the following diagram commutes for all 1-cells \(f \in \mathbb{B}(x,y)\):
  \[
    \begin{mytikzcd}[ampersand replacement=\&, row sep=small, column sep=small]
      Fx \&\& Fy \&\& Fx \&\& Fy \\
      \&\&\& {=} \\
      Gx \&\& Gy \&\& Gx \&\& Gy
      \arrow["Ff", from=1-1, to=1-3]
      \arrow[""{name=0, anchor=center, inner sep=0}, "{\alpha_x}", curve={height=-12pt}, from=1-1, to=3-1]
      \arrow[""{name=1, anchor=center, inner sep=0}, "{\beta_x}"', curve={height=12pt}, from=1-1, to=3-1]
      \arrow["{\alpha_f}", shift left=4, shorten <=14pt, shorten >=14pt, Rightarrow, from=1-3, to=3-1]
      \arrow["{\alpha_y}", from=1-3, to=3-3]
      \arrow["Ff", from=1-5, to=1-7]
      \arrow["{\beta_x}"', from=1-5, to=3-5]
      \arrow["{\beta_f}", shift right=6, shorten <=14pt, shorten >=14pt, Rightarrow, from=1-7, to=3-5]
      \arrow[""{name=2, anchor=center, inner sep=0}, "{\beta_y}"', curve={height=12pt}, from=1-7, to=3-7]
      \arrow[""{name=3, anchor=center, inner sep=0}, "{\alpha_y}", curve={height=-12pt}, from=1-7, to=3-7]
      \arrow["Gf"', from=3-1, to=3-3]
      \arrow["Gf"', from=3-5, to=3-7]
      \arrow["{\raisebox{-0.7em}{\(\Upsilon_x\)}}", shorten <=4pt, shorten >=4pt, Rightarrow, from=0, to=1]
      \arrow["{\raisebox{-0.7em}{\(\Upsilon_y\)}}", shorten <=4pt, shorten >=4pt, Rightarrow, from=3, to=2]
    \end{mytikzcd}
  \]
\end{definition}

In fact, bicategories, lax functors, lax natural transformations, and modifications between them
form a \emph{tricategory}.
We shall not give the formal definition here,
but see~\cite{gordon95:coher,gurski2006algebraic}
and~\cite[Chapter~11]{johnson-yau2021:TwoDimensionalCategories}.

\section{Monads and adjunctions}\label{sec:monads-adjunctions}

\begin{definition}\label{def:monad}
  A \emph{monad} on a category \(\cat{C}\) consists of an endofunctor \(T\) on \(\cat{C}\)
  and two natural transformations
  \(\mu \from T^2 \nt T\) and \(\eta \from \Id_{\cat{C}} \nt T\),
  called the \emph{multiplication} and \emph{unit} of \(T\),
  satisfying \emph{associativity} and \emph{unitality} axioms:
  \[
    \begin{mytikzcd}[ampersand replacement=\&]
      {T^3} \& { T^2} \& {T^2} \& T \& T \\
      {T^2} \& T \&\& T
      \arrow["T\mu", rightarrow, from=1-1, to=1-2]
      \arrow["\mu", rightarrow, from=1-2, to=2-2]
      \arrow["{\mu T}"', rightarrow, from=1-1, to=2-1]
      \arrow["\mu"', rightarrow, from=2-1, to=2-2]
      \arrow[rightarrow, no head, from=1-4, to=2-4]
      \arrow["{\eta T}"', rightarrow, from=1-4, to=1-3]
      \arrow["\mu"', rightarrow, from=1-3, to=2-4]
      \arrow["{T \eta}", rightarrow, from=1-4, to=1-5]
      \arrow["\mu", rightarrow, from=1-5, to=2-4]
    \end{mytikzcd}
  \]
\end{definition}

\begin{definition}\label{def:adjunction}
  An \emph{adjunction}
  consists of a pair of functors \(\stdadj\),
  together with two natural transformations,
  the \emph{unit} \(\eta\from \Id_{\cat{C}} \nt UF\)
  and the \emph{counit} \(\varepsilon\from FU \nt \Id_{\cat{D}}\),
  satisfying the \emph{snake} or \emph{triangle} identities;
  see for example~\cite[Definition~4.2.5]{category-theory-in-context_riehl}.
\end{definition}

Alternatively,
one could define an adjunction between \(F\from \cat{C}\to \cat{D}\) and \(G\from \cat{D}\to \cat{C}\)
to require the existence of a natural isomorphism
\[
  \cat{D}(F(\blank), \bblank) \cong \cat{C}(\blank, U(\bblank)),
\]
from which one recovers the unit and counit, see~\cite[Proposition~4.2.6]{category-theory-in-context_riehl}.
Note that, a priori,
different natural isomorphisms lead to different adjunctions between the same functors.
However, fixing for example \(G\),
then for two left adjoints \(F\) and \(F'\)
there exists a unique isomorphism \(\Theta\from F \nt F'\)
that commutes with the respective units and counits;
see~\cite[Proposition~4.4.1]{category-theory-in-context_riehl}.

\begin{example}\label{ex:interplay-monad-adjunction}
  If \(\stdadj\) is an adjunction with unit \(\eta\) and counit \(\varepsilon\),
  then \(UF \from \cat{C} \to \cat{C}\) is a monad;
  the multiplication \(\mu^T\from UFUF \nt UF\) is given by \(U \varepsilon F\),
  and the unit \(\eta^T\from \Id_{\cat{C}} \nt UF\) is equal to \(\eta\).
\end{example}

\begin{definition}\label{def:algebras-over-monad}
  Given a monad \((T, \mu, \eta)\) on a category \(\cat{C}\),
  a \emph{\(T\)\hyp{}algebra} consists of an object \(x \in \cat{C}\) and a morphism \(\alpha \from Tx \to x\),
  satisfying
  \[
    \alpha \circ T\alpha = \alpha \circ \mu_x
    \qquad\qquad\text{and}\qquad\qquad
    \alpha \circ \eta_x = \id_x.
  \]

  Given two \(T\)\hyp{}algebras \((x, \alpha)\) and \((y, \beta)\),
  a \emph{morphism} between them consists of a morphism \(f\from x \to y\) in \(\cat{C}\),
  such that \(\beta \circ Tf = f \circ \alpha\).
\end{definition}

\begin{remark}\label{rmk:em-is-algebras}
  For any monad \(T\), the \(T\)\hyp{}algebras and their morphisms form a category: the \emph{\EilenbergMoore{} category of \(T\)}.%
  \marginnote{%
    Our notation for the Eilenberg–Moore category
    already ascribes a certain ``terminal'' quality to it;
    see \cref{rmk:em-kl-terminal-initial}.%
  }
  We shall denote it by \(\cat{C}^T\).

  The Eilenberg–Moore category of \(T\) is also often called the \emph{category of \(T\)\hyp{}algebras} or,
  following for example~\cite{Bruguieres2007},
  the \emph{category of modules} over \(T\).
  We use all three terminologies interchangeably.
\end{remark}

A monad is intimately connected to its Eilenberg–Moore category.

\begin{example}\label{ex:eilenberg-moore-adjunctions}
  There is a \(2\)\hyp{}category \(\BiCat{Mon}(\BiCat{Cat})\) of monads in \(\BiCat{Cat}\),~\cite[\S~1]{Street1972}.
  The inclusion 2-functor maps a category to its identity monad:
  \[
    \BiCat{Cat} \to \BiCat{Mon}(\BiCat{Cat}),\qquad \cat{C} \mapsto (\Id_{\cat{C}}, \id_{\Id_{\cat{C}}}, \id_{\Id_{\cat{C}}}).
  \]
  By assumption, \(\BiCat{Cat}\) \emph{admits the construction of algebras}:
  there exists a right adjoint to the above functor:
  \[
    \BiCat{Mon}(\BiCat{Cat}) \to \BiCat{Cat},
    \qquad
    (T\from \cat{C}\to \cat{C}, \mu, \eta) \mapsto \cat{C}^{T},
  \]
  where \(\cat{C}^{T} \in \BiCat{Cat}\) is the \EilenbergMoore{} category of \(T\).
  Using the previous 2-adjunction,
  one proves that to every monad \((T, \mu, \eta)\) on \(\cat{C}\)
  there exist an \emph{\EilenbergMoore{} adjunction}
  \(F^T \from \cat{C} \to \cat{C}^T\)
  and
  \(U^T \from \cat{C}^T \to \cat{C}\),
  such that
  \[
    T = U^T F^T,\qquad
    \mu = F^T \varepsilon U^T,\qquad
    \eta = \eta,
  \]
  where
  \(\eta \from 1_{\cat{C}} \nt U^T F^T\) and \(\varepsilon \from F^T U^T \nt 1_{\cat{C}^T}\)
  are the unit and counit of the \EilenbergMoore{} adjunction.
  We shall call \(F^T\) and \(U^T\) the \emph{free} and \emph{forgetful} functor \emph{associated to \(T\)}, respectively.
\end{example}

For the following definition, we follow~\cite{Bruguieres2007,Turaev2017}.

\begin{definition}\label{def:morphism-of-monads}
  Suppose that \(T\) and \(S\) are two monads on the category \(\cat{C}\).
  A \emph{morphism of monads} between \(T\) and \(S\) is a natural transformation \(f \from T \nt S\),
  such that the following diagrams commute
  \begin{equation}\label[diagram]{eq:morphism-of-monads}
    \begin{mytikzcd}[ampersand replacement=\&]
      \Id \& T \&\&\& TT \& TS \& SS \\
      \& S \&\&\& T \&\& S
      \arrow["{\eta^T}", from=1-1, to=1-2]
      \arrow["{\eta^S}"', from=1-1, to=2-2]
      \arrow["f", from=1-2, to=2-2]
      \arrow["Tf", from=1-5, to=1-6]
      \arrow["{\mu^T}"', from=1-5, to=2-5]
      \arrow["fS", from=1-6, to=1-7]
      \arrow["{\mu^S}", from=1-7, to=2-7]
      \arrow["f"', from=2-5, to=2-7]
    \end{mytikzcd}
  \end{equation}
\end{definition}

\begin{remark}
  The terminology of \cref{def:morphism-of-monads} is slightly non-standard.
  What we call a morphism of monads is often called a \emph{oplax} (or colax) monad morphism,
  see for example~\cite[\S~1]{Street1972}.
\end{remark}

\begin{remark}
  One can define a monad in any bicategory \(\mathbb{B}\) by considering \((C, t, \eta, \mu)\),
  where \(C \in\mathbb{B}\) is an object,
  \(t\from C \to C\) is a 1-cell,
  and \(\eta\from \Id_C \nt t\) and \(\mu\from tt \nt t\) are 2-cells,
  satisfying relations analogous to \cref{def:monad}.
  An oplax morphism of monads from \((C, t, \eta^t, \mu^t)\) to \((D, s, \eta^s, \mu^s)\)
  then consists of a 1-cell \(u\from C \to D\) and a 2-cell \(\phi\from ut \nt su\),
  subject to identities reminiscent of \cref{eq:morphism-of-monads}.
  A \emph{lax morphism of monads} involves a 1-cell \(u\from C \to D\) and a 2-cell \(\phi\from su \nt ut\),
  satisfying similar properties.
\end{remark}

The following example sheds some additional light on this terminology.

\begin{example}\label{ex:alternative-monad}
  Monads can alternatively be defined as
  lax functors—in the sense of \cref{def:bifunctor}—from
  the terminal 2-category \(\terminal\) to \(\BiCat{Cat}\).
  Unravelling this definition,
  a lax functor \(\mathfrak{T} \from \terminal \to \BiCat{Cat}\) consists of:
  \begin{itemize}
    \item an object assignment \(\mathfrak{T}\from \Ob\terminal \to \Ob\BiCat{Cat}\),
    sending the unique object \(\ast\) to a category \(\cat{C}\);
    \item a functor \(\mathfrak{T}(*,*)\from \terminal(*,*) \to \BiCat{Cat}(\cat{C},\cat{C})\)
    from the terminal 1-category \(\terminal(*,*)\) to the category of endofunctors on \(\cat{C}\),
    sending the unique 1-morphism \(\id_{*}\from * \to *\) to \(T\from \cat{C}\to \cat{C}\)
    and the unique 2-morphism \(1_{\id_{*}}\from \id_{*} \nt \id_{*}\) to the identity natural transformation \(T \nt T\);
    \item a 2-cell\marginnote{%
      Recall that \(\otimes\) is the horizontal composition in \(\mathbb{B}\).
    }
    \(\mathfrak{T}_2\from \mathfrak{T}\id_{*} \otimes \mathfrak{T}\id_{*} \nt \mathfrak{T}\id_{*}\),
    which we write as \(\mu\from T T \nt T\); and
    \item a 2-cell
    \(\mathfrak{T}_0\from 1_{\mathfrak{T}(*)} \nt \mathfrak{T}\id_{*}\)
    that we write as \(\eta \from \Id_{\cat{C}} \nt T\).
  \end{itemize}
  The properties of \cref{def:bifunctor} for \(\mathfrak{T}_2\) and \(\mathfrak{T}_0\)
  translate to the associativity and unitality properties of \(\mu\) and \(\eta\).
  In this setting,
  a morphism of monads becomes an oplax transformation in the sense of \cref{def:oplax-transformation}.
\end{example}

\begin{example}\label{ex:kleisli-category}
  A monad \(T\) on \(\cat{C}\) has another canonical category associated to it:
  its \emph{Kleisli category} \(\cat{C}_T\).
  On objects, it is given by \(\mathrm{Ob}(\cat{C}_T) \defeq \mathrm{Ob}(\cat{C})\),
  and for \(x, y \in \cat{C}_T\) we have \(\cat{C}_T(x,y) \defeq \cat{C}(x,Ty)\).
  Composition is defined by
  \begin{align*}
    \circ \from \cat{C}_T(y, z) \times \cat{C}_T(x, y) &\to \cat{C}_T(x, z) \\
    (g, f) &\mapsto \big( x \xrightarrow{\ \ f\ \ } Ty \xrightarrow{\ Tg\ } T^2z \xrightarrow{\ \ \mu_z\ \ } Tz \big).
  \end{align*}
\end{example}

\begin{proposition}\label{prop:Kleisli-adjunction}
  Let \(T\) be a monad on a category \(\cat{C}\).
  There exists an adjunction
  \[
    \begin{mytikzcd}[ampersand replacement=\&]
      {\cat{C}} \& {\cat{C}_T}
      \arrow[""{name=0, anchor=center, inner sep=0}, "{F_T}", shift left=2, from=1-1, to=1-2]
      \arrow[""{name=1, anchor=center, inner sep=0}, "{U_T}", shift left=2, from=1-2, to=1-1]
      \arrow["\dashv"{anchor=center, rotate=-90}, draw=none, from=0, to=1]
    \end{mytikzcd}
  \]
  where \(F_T\) is identity on objects and sends \(f \in \cat{C}(x,y)\) to \(\eta_{y} \circ f\),
  and \(U_T\) sends \(x\) to \(Tx\) and \(f \in \cat{C}(x, y)\) to \(\mu_{y} \circ Tf\).
\end{proposition}

\subsection{Comparison functors}\label{sec:comp-funs}

\textsc{Given a monad \(T\) coming from} the adjunction \stdadj,
we might ask how much the functors \(F\) and \(U\) ``differ''
from the free and forgetful functors \(F^T \from \cat{C} \to \cat{C}^T\) and \(U^T \from \cat{C}^T \to \cat{C}\) of \(T\).
Roughly summarised we are interested in the following:
\[
  \begin{mytikzcd}[ampersand replacement=\&]
    \cat{D} \arrow[rrdd, "U", bend left] \& \& \& \& \cat{C}^T \arrow[llll, "\text{``compare''}"', no head, dashed, bend right] \arrow[lldd, "U^T", bend left] \\
    \&  \&                                                                        \&  \&                                                                                          \\
    \&  \& \cat{C} \arrow[lluu, "F", bend left] \arrow[rruu, "F^T", bend left] \&  \&
  \end{mytikzcd}
\]

\begin{proposition}[{\cite[Theorem~2]{frei69:some}}]\label{prop:morphism-of-monads-gets-kleisli}
  Let \(T\) and \(S\) be monads on a category \(\cat{C}\).
  Then there exists a bijection between monad morphisms from \(T\) to \(S\)
  and functors \(F\from \cat{C}^S \to \cat{C}^T\) between their categories of algebras,
  such that \(U^T F = U^S\).
\end{proposition}

Note that, if we were to define morphisms of monads as lax transformations,
\cref{prop:morphism-of-monads-gets-kleisli} would yield a covariant assignment.

\begin{lemma}[{\cite[Theorem~3]{Street1972}}]\label{lem:comparison-functor}
  Let \(\stdadj\) be an adjunction
  and set \(T \defeq UF\).
  There exists a unique functor
  \(K^T \from \cat{D} \to \cat{C}^T\)
  satisfying \(K^T F = F^T\) and \(U^T K^T = U\).
  On objects it is given by \(K^T d = (Ud, U\varepsilon_{d})\), for all \(d \in \cat{D}\).
\end{lemma}

We call the unique functor from \cref{lem:comparison-functor}
the \emph{comparison functor}.
An adjunction is called \emph{monadic} if its comparison functor is an equivalence.

There exists an analogous version of \cref{lem:comparison-functor} for the Kleisli category of a monad \(T\).
The comparison functor in this situation is given by
\[
  K_{T}\from \cat{C}_T \to \cat{D}, \qquad x \mapsto Fx, \qquad f \in \cat{D}(x,Ty) \mapsto \varepsilon_{Fy} \circ Ff.
\]

\begin{remark}\label{rmk:em-kl-terminal-initial}
  In fact, the Eilenberg--Moore and the Kleisli category are the terminal and initial objects in
  the suitably defined category of adjunctions \(\stdadj\) producing the monad \(T\).
  Thus, the diagram
  \[
    \begin{mytikzcd}[ampersand replacement=\&]
      {\cat{C}_T} \&\& {\cat{D}} \&\& {\cat{C}^T} \\
      \\
      \&\& {\cat{C}}
      \arrow["{K_T}", from=1-1, to=1-3]
      \arrow[""{name=0, anchor=center, inner sep=0}, "{U_T}"{pos=0.4}, shift left, from=1-1, to=3-3]
      \arrow["{K^T}", from=1-3, to=1-5]
      \arrow[""{name=1, anchor=center, inner sep=0}, "U", shift left=2, from=1-3, to=3-3]
      \arrow[""{name=2, anchor=center, inner sep=0}, "{U^T}", shift left=3, from=1-5, to=3-3]
      \arrow[""{name=3, anchor=center, inner sep=0}, "{F_T}", shift left=3, from=3-3, to=1-1]
      \arrow[""{name=4, anchor=center, inner sep=0}, "F", shift left=2, from=3-3, to=1-3]
      \arrow[""{name=5, anchor=center, inner sep=0}, "{F^T}"{pos=0.6}, shift left, from=3-3, to=1-5]
      \arrow["\dashv"{anchor=center}, draw=none, from=4, to=1]
      \arrow["\dashv"{anchor=center, rotate=49}, draw=none, from=3, to=0]
      \arrow["\dashv"{anchor=center, rotate=-50}, draw=none, from=5, to=2]
    \end{mytikzcd}
  \]
  commutes,
  and its commutativity characterises \(K^T\) and \(K_T\) completely.
\end{remark}

The composite \(K^T \circ K_T\) is fully faithful,
its full image consisting of \emph{free modules}.
We will denote the canonical inclusion of \(\cat{C}_T\) into \(\cat{C}^T\) by
\begin{equation}\label{denotedbyiota}
  \iota\from \cat{C}_T \longhookrightarrow \cat{C}^T, \qquad x \mapsto (Tx, \mu_x), \qquad f \in \cat{C}(x, Ty) \mapsto \mu_y \circ Tf.
\end{equation}

\begin{example}\label{ex:comonads}
  A \emph{comonad} on a category \(\cat{C}\) consists of
  an endofunctor \(S\) on \(\cat{C}\) together with natural transformations \(\Delta\from S \nt S^2\) and \(\varepsilon\from S \nt \Id_{\cat{C}}\),
  satisfying axioms dual to those of \cref{def:monad}.

  For a comonad \(S\) on \(\cat{C}\)
  an \emph{\(S\)\hyp{}comodule} or \emph{\(S\)\hyp{}coalgebra} is an object \(x \in \cat{C}\) together with a morphism \(x \to S x\) satisfying axioms analogous to \cref{def:algebras-over-monad}.
  We will denote the category of \(S\)\hyp{}comodules by \(\cat{C}^{S}\) and the Kleisli category of \(S\) by \(\cat{C}_{S}\).
  They have analogously defined adjunctions
  \[
    \begin{mytikzcd}[ampersand replacement=\&]
      {\cat{C}_{S}} \& {\cat{C}} \& {\text{and}} \& {\cat{C}^{S}} \& {\cat{C}.}
      \arrow[""{name=0, anchor=center, inner sep=0}, "{F_{S}}", shift left=2, from=1-1, to=1-2]
      \arrow[""{name=1, anchor=center, inner sep=0}, "{U_{S}}", shift left=2, from=1-2, to=1-1]
      \arrow[""{name=2, anchor=center, inner sep=0}, "{F^{S}}", shift left=2, from=1-4, to=1-5]
      \arrow[""{name=3, anchor=center, inner sep=0}, "{U^{S}}", shift left=2, from=1-5, to=1-4]
      \arrow["\dashv"{anchor=center, rotate=-90}, draw=none, from=0, to=1]
      \arrow["\dashv"{anchor=center, rotate=-90}, draw=none, from=2, to=3]
    \end{mytikzcd}
  \]

  The Eilenberg--Moore category and the Kleisli category for the comonad \(S \defeq FU\) can be characterised as
  a terminal and initial object, respectively.
  This yields functors \(K_{S}\) and \(K^{S}\) such that the following diagram commutes
  \[
    \begin{mytikzcd}[ampersand replacement=\&]
      \&\& {\cat{D}} \\
      \\
      {\cat{C}_{S}} \&\& {\cat{C}} \&\& {\cat{C}^{S}}
      \arrow[""{name=0, anchor=center, inner sep=0}, "{{U_{S}}}"{pos=0.6}, shift left, from=1-3, to=3-1]
      \arrow[""{name=1, anchor=center, inner sep=0}, "U", shift left=2, from=1-3, to=3-3]
      \arrow[""{name=2, anchor=center, inner sep=0}, "{{U^{S}}}", shift left=3, from=1-3, to=3-5]
      \arrow[""{name=3, anchor=center, inner sep=0}, "{{F_{S}}}", shift left=3, from=3-1, to=1-3]
      \arrow["{{K_{S}}}"', from=3-1, to=3-3]
      \arrow[""{name=4, anchor=center, inner sep=0}, "F", shift left=2, from=3-3, to=1-3]
      \arrow["{{K^{S}}}"', from=3-3, to=3-5]
      \arrow[""{name=5, anchor=center, inner sep=0}, "{{F^{S}}}"{pos=0.4}, shift left, from=3-5, to=1-3]
      \arrow["\dashv"{anchor=center, rotate=-48}, draw=none, from=3, to=0]
      \arrow["\dashv"{anchor=center}, draw=none, from=4, to=1]
      \arrow["\dashv"{anchor=center, rotate=49}, draw=none, from=5, to=2]
    \end{mytikzcd}
  \]

  Similarly to \cref{ex:interplay-monad-adjunction},
  given an adjunction \(\stdadj\), we obtain a comonad \(FU\) on \(\cat{D}\).
\end{example}

It may be that a monad \(T\) on \(\cat{C}\) has a right adjoint \(G\)%
—in this case, one can automatically equip \(G\) with a canonical comonad structure.

\begin{proposition}[{\cite[Theorem~V.8.2]{mac92:sheav}}]\label{pyramidscheme}
  If\, \(T\) is a monad on a category \(\cat{C}\) and \(G\) is a right adjoint to \(T\),
  then \(G\) is a comonad and there is a canonical isomorphism between the Eilenberg–Moore categories: \(\cat{C}^T \cong \cat{C}^G\).
\end{proposition}

\begin{proposition}[{\cite[Theorem~3]{kleiner90:adjoin-kleis}}]\label{einerKleinerKleisliÄquivalenz}
  If \(T\) is a monad on \(\cat{C}\) and \(S\) is a left adjoint to \(T\),
  then \(S\) is a comonad and there is a canonical isomorphism \(\cat{C}_T \cong \cat{C}_S\).
\end{proposition}

In general, it is \emph{not} true that \(\cat{C}^T \cong \cat{C}^S\) in the setting of \cref{einerKleinerKleisliÄquivalenz}.

\subsection{Distributive laws}\label{sec:distributive-laws}

\textsc{Beck's theory of distributive laws}
concerns itself with the question when, given a monad \(T\) on a category \(\cat{C}\),
a functor \(S\from \cat{C} \to \cat{C}\) \emph{lifts} to a functor on the Eilenberg–Moore category of \(T\),~\cite{Beck1969}.
That is, there is some \(\tilde{S}\from \cat{C}^T\to \cat{C}^T\), that satisfies the lifting condition
\[
  \begin{mytikzcd}[ampersand replacement=\&]
    \& {\cat{C}^T} \\
    {\cat{C}^T} \& {\cat{C}}
    \arrow["{U^T}", from=1-2, to=2-2]
    \arrow["{\tilde{S}}", dashed, from=2-1, to=1-2]
    \arrow["{S \circ U^T}"', from=2-1, to=2-2]
  \end{mytikzcd}
\]
If \(S\) is a monad this in turn is equivalent to  \(ST\) being a monad itself.
Distributive laws thus yield a witness for the composability of two monads.
Street further developed the theory of distributive laws intrinsic to certain well-behaved bicategories, see~\cite{Street1972}.

\begin{definition}\label{def:distributive-law}
  Let \(T\) and \(S\) be two monad on a category \(\cat{C}\).
  A \emph{distributive law} of \(T\) over \(S\) consists of a natural transformation \(\Omega\from TS \nt ST\)
  such that the following diagrams commute:
  \[
    \begin{mytikzcd}[ampersand replacement=\&]
      TSS \& STS \& SST \& {T } \& TS \\
      TS \&\& ST \&\& ST
      \arrow["{T \mu^S}"', from=1-1, to=2-1]
      \arrow["\Omega"', from=2-1, to=2-3]
      \arrow["{\Omega S}", from=1-1, to=1-2]
      \arrow["{S \Omega}", from=1-2, to=1-3]
      \arrow["{\mu^S T}", from=1-3, to=2-3]
      \arrow["{T \eta^S}", from=1-4, to=1-5]
      \arrow["\Omega", from=1-5, to=2-5]
      \arrow["{\eta^S T}"', from=1-4, to=2-5]
    \end{mytikzcd}
  \]
  \[
    \begin{mytikzcd}[ampersand replacement=\&]
      TTS \& TST \& STT \& S \& TS \\
      TS \&\& ST \&\& ST
      \arrow["{T \Omega}", from=1-1, to=1-2]
      \arrow["{\Omega T}", from=1-2, to=1-3]
      \arrow["{S \mu^T}", from=1-3, to=2-3]
      \arrow["{\mu^T S}"', from=1-1, to=2-1]
      \arrow["\Omega"', from=2-1, to=2-3]
      \arrow["{S \eta^T}"', from=1-4, to=2-5]
      \arrow["{\eta^T S}", from=1-4, to=1-5]
      \arrow["\Omega", from=1-5, to=2-5]
    \end{mytikzcd}
  \]
\end{definition}

One analogously defines a distributive law between two comonads,
and \emph{mixed distributive laws} between a monad and a comonad.

\begin{theorem}[{\cite{Beck1969}}]
  Let \(S\) and \(T\) be monads on a category \(\cat{C}\).
  Then the following are equivalent:
  \begin{thmlist}[noitemsep]
    \item distributive laws of \(T\) over \(S\);
    \item monad structures \(S \circ_{\Omega} T \defeq (ST, \mu, \eta)\) on \(ST\),
    such that \(S \eta^T\) and \(\eta^S T\) are morphisms of monads
    and \(1_{ST} = \mu \circ S \eta^T \eta^T T\); and
    \item lifts of \(S\) to the Eilenberg–Moore category of \(T\),
    such that \(S\) is a monad on \(\cat{C}^T\).
  \end{thmlist}
\end{theorem}
\begin{proof}[Sketch of proof]
  We outline the main constructions.

  Given a distributive law \(\Omega\from TS \nt ST\),
  for two monads on \(\cat{C}\),
  the monad structure on \(ST\) is given by
  \[
    \mu\from STST \xrightarrow{S \Omega T} SS TT \xrightarrow{\mu^S \mu^T} ST
    \qquad\text{and}\qquad
    \eta\from \Id_{\cat{C}} \xrightarrow{\eta^T} T \xrightarrow{\eta^S T} ST.
  \]
  The lift \(\tilde{S}\) of \(S\) to \(\cat{C}^T\) is defined by
  \(\tilde{S}(x, \nabla_x) \defeq (Sx, TSx \xrightarrow{\Omega_x} STx \xrightarrow{S \nabla_x} Sx)\).
  Its multiplication and unit is given by those of \(S\).

  Given a monad structure \((ST, \mu, \eta)\),
  one defines a distributive law by
  \[
    TS \xrightarrow{\eta^S TS \eta^T} STST \xrightarrow{\ \mu\ } ST.
  \]

  Finally, if \(\tilde{S}\) is a lift of \(S\),
  we obtain a distributive law via
  \[
    TS \xrightarrow{TS \eta^{T}} TS U^T F^T = U^T F^T U^T \tilde{S} F^T \xrightarrow{U^{T} \varepsilon^T \tilde{S} F^T} U^T \tilde{S} F^T = S T.\qedhere
  \]
\end{proof}

\section{String diagrams}\label{sec:graphical-calculus-introduction}

\textsc{Recall from \cref{ex:cat-is-a-bicategory} that}
\(\BiCat{Cat}\) is the 2-category of all (small) categories, functors, and natural transformations.
We use juxtaposition for the horizontal composition of \(\BiCat{Cat}\), or, in case we want to emphasise the direction of composition,
\[
  \blank \odot \bblank \from \BiCat{Cat}(\cat{B},\cat{C}) \times \BiCat{Cat}(\cat{A},\cat{B}) \to \BiCat{Cat}(\cat{A},\cat{C}),
  \qquad \qquad
  \cat{A},\cat{B},\cat{C} \in \BiCat{Cat}.
\]

String diagrams will serve as an important tool for doing computations.\marginnote{%
  More generally, our applications of the graphical calculus can be formulated in any bicategory
  that admits the construction of algebras in the sense of \cref{ex:eilenberg-moore-adjunctions}.
  For ease of presentation we shall stick with \(\BiCat{Cat}\).
}
In the case of \(\BiCat{Cat}\),
a \emph{string diagram} consists of
regions labelled with categories,
strings labelled with functors,
and vertices between the strings labelled with natural transformations.
If two string diagrams can be transformed into each other%
—that is, if they are isotopic—%
then the natural transformations they represent are equal.
A more detailed description is given in~\cite{joyal91,Selinger2011}.
Our convention is to read diagrams from bottom to top and right to left.
Horizontal and vertical composition are given by horizontal and vertical gluing of  diagrams, respectively.
Identity natural transformations are given by unlabelled vertices.
The edge for the identity functor of a category will not be drawn.
If the involved categories are clear from the context, we omit writing them explicitly.
\Cref{fig:explanation-string-diagrams} details our conventions.
\begin{figure}[htbp]
  \centering
  \tikzfig{explanation-string-diagrams}\vspace{-0.2cm}% XXX
  \scaption[1.2cm]{%
    Basic string diagrammatic conventions.%
    \label{fig:explanation-string-diagrams}%
  }
\end{figure}

\begin{example}\label{ex:graphical-monads}
  Let \(T\) be a monad on \(\cat{C}\).
  We represent the multiplication and unit of \(T\) in terms of string diagrams:
  \[
    \tikzfig{monads-multiplication-and-unit}
  \]
  Their associativity and unitality then equate to
  \[
    \tikzfig{monads-assoc-and-unit}
  \]

  We depict the unit and counit of the \EilenbergMoore{} adjunction as
  \[
    \tikzfig{adjunction-unit-counit}
  \]
  The defining equations of adjunctions translate to the snake equations
  \[
    \tikzfig{adjunction-snake-identities}
  \]
\end{example}

\begin{remark}\label{rmk:diagrammatic-algebras}
  Given a monad \(T\), its category of algebras can be incorporated into the graphical calculus, see~\cite{Willerton2008}.
  Define the natural transformation
  \begin{equation}\label{eq:action-as-nat-trafo}
    \nabla \from T U^T = U^T F^T U^T \xrightarrow{\ U^T\varepsilon\ } U^T.
  \end{equation}
  Now, one may write
  \[
    \tikzfig{action-over-monad}
  \]
  We think of \(\nabla \from T U^T \to U^T\) as
  an action of \(T\) on \(U^T\) due to the identities
  \[
    \tikzfig{axioms-of-action-over-monad}
  \]

  By abuse of notation, we shall often simply speak of an object \(x \in \cat{C}^T\),
  leaving the action \(\nabla_x\) implicit.
\end{remark}

\section{Monoidal and module categories}\label{sec:mono-module-categ}

\textsc{For a more extensive account}
concerning the notions of monoidal and module categories,
we refer the reader to~\cite[Chapter~2]{Etingof2015}.

\begin{definition}\label{def:monoidal-category}
  A \emph{monoidal category} consists of a category \(\cat{C}\)
  together with a \emph{tensor product} functor \(\otimes \from \cat{C} \times \cat{C} \to \cat{C}\),
  a \emph{unit} \(1 \in \cat{C}\),
  and for all \(x, y, z \in \cat{C}\), natural isomorphisms
  \begin{align*}
    \alpha_{x, y, z} \from (x \otimes y) \otimes z \Iso x \otimes (y \otimes z),
    \quad\ \ %
    \lambda_{x} \from 1 \otimes x \Iso x,
    \quad\ \ %
    \rho_{x} \from x \otimes 1 \Iso x,
  \end{align*}
  satisfying \emph{coherence} axioms;
  see~\cite[Definition~2.1.1]{Etingof2015}.

  A monoidal category is called \emph{strict} if \(\alpha\), \(\lambda\), and \(\rho\) are identities.
\end{definition}

\begin{example}\label{ex:monoidal-rev}
  Given a monoidal category \((\cat{C}, \otimes, 1)\),
  we denote by \(\cat{C}^{\tensorop}\) the monoidal category \((\cat{C}, \otimes^{\op}, 1)\),
  where for all \(x, y \in \cat{C}\) one sets \(x \otimes^{\op} y \defeq y \otimes x\).
\end{example}

Recall \emph{Sweedler notation}:
if \(B\) is a bialgebra in \(\kVect\) with comultiplication \(\Delta\from B \to B \kotimes B\),
then for \(b \in B\) we write \(b_{(1)} \otimes b_{(2)} \defeq \Delta(b) \in B \kotimes B\).
For example, in this notation coassociativity translates to
\[
  b_{(1)} \otimes {(b_{(2)})}_{(1)} \otimes {(b_{(2)})}_{(2)}
  = b_{(1)} \otimes b_{(2)} \otimes b_{(3)}
  = {(b_{(1)})}_{(1)} \otimes {(b_{(1)})}_{(2)} \otimes b_{(2)}.
\]
This naturally extends to left and right coactions on \(B\),
which we respectively denote by \(x \mapsto x_{(-1)} \otimes x_{(0)}\) and \(x \mapsto x_{(0)} \otimes x_{(1)}\).

\begin{example}[{\cite{yetter90:quant}}]\label{ex:yd-modules-def-and-monoidal}
  Let \(H \in \kVect\) be a bialgebra
  The category \(\YDM{H}\) of (left-left) \emph{\YetterDrinfeld{} modules} is defined as the category
  of those \(Y \in \kVect\) that are left \(H\)\hyp{}modules with action \(\lact\)
  and also left \(H\)\hyp{}comodules with coaction \(y \mapsto y_{(-1)} \otimes y_{(0)}\),
  such that the following compatibility condition holds for all \(h \in H\) and \(y \in Y\):
  \begin{equation}\label{eq:yd-condition-bialgebra}
    h_{(1)}y_{(-1)} \otimes (h_{(2)} \lact y_{(0)}) = {(h_{(1)} \lact y)}_{(-1)} h_{(2)} \otimes {(h_{(1)} \lact y)}_{(0)}
  \end{equation}

  Then \(\YDM{H}\) is monoidal: given \YetterDrinfeld{} modules \(Y\) and \(Z\),
  their tensor product \(Y \kotimes Z\) can be equipped with the structure of a Yetter–Drinfeld module.
  For all \(h \in H\), \(y \in Y\), and \(z \in Z\), define
  \[
    h \lact (y \otimes z) \defeq (h_{(1)} \lact y) \otimes (h_{(2)} \lact z)
    \quad\text{and}\quad
    y \otimes z \mapsto y_{(-1)}z_{(-1)} \otimes y_{(0)} \otimes z_{(0)}.
  \]

  In fact, if \(H\) is a Hopf algebra with invertible antipode \(S\)%
  —for example, when \(H\) is finite\hyp{}dimensional—%
  then \cref{eq:yd-condition-bialgebra} is equivalent to
  \begin{equation}\label{eq:yd-condition}
    {(h \lact y)}_{(-1)} \otimes {(h \lact y)}_{(0)} = h_{(1)}y_{(-1)}S(h_{(3)}) \otimes (h_{(2)} \lact y_{(0)}),
  \end{equation}
  see~\cite[Proposition~10.6.16]{montgomery93:hopf}.
\end{example}

\begin{definition}\label{def:monoidal-functors}
  A \emph{lax monoidal functor} between monoidal categories \(\cat{C}\)
  and \(\cat{D}\)
  consists of a functor \(F\from \cat{C} \to \cat{D}\)
  together with a morphism \(F_0 \from 1 \to F1\)
  and a natural transformation \(F_2 \from F(\blank) \otimes F(\bblank) \nt F(\blank \otimes \bblank)\),
  satisfying \emph{associativity} and \emph{unitality} conditions, see~\cite[Definition~2.4.1]{Etingof2015}.

  Analogously, an \emph{oplax monoidal functor} is a functor \(G\from \cat{C} \to \cat{D}\)
  together with morphisms \(G_0\from G1 \to 1\) and \(G_2 \from G(\blank \otimes \bblank) \nt G(\blank) \otimes G(\bblank)\),
  making \(G^{\op}\from \cat{C}^{\op} \to \cat{D}^{\op}\) lax monoidal.
  We call \(F\) \emph{strong monoidal}
  if \(F_0\) and \(F_2\) and invertible,
  and \emph{strict monoidal} if they are identities.
\end{definition}

Observe that an oplax monoidal functor with invertible coherence morphisms canonically defines a strong monoidal functor.

\begin{remark}\label{rmk:whynotthree}
  Let \(T\) be an oplax monoidal endofunctor on a monoidal category \(\cat{C}\).
  By coassociativity of \(T_2\),
  there exists a well-defined natural transformation
  \[
    T_3\from T(\blank \otimes \bblank \otimes \bbblank) \nt T(\blank) \otimes T(\bblank) \otimes T(\bbblank),
  \]
  which, for all \(x,y,z \in \cat{C}\), has components
  \[
    T_{3;x,y,z}
    \defeq (T_{2;x,y} \otimes Tz) \circ T_{2; x \otimes y, z}
    = (Tx \otimes T_{2; y, z}) \circ T_{2; x, y \otimes z}.
  \]
\end{remark}

\begin{definition}\label{def:monoidal-nat-trafo}
  Let \(\cat{C}\) and \(\cat{D}\) be two monoidal categories and suppose that \(F,G \colon \cat{C} \to \cat{D}\) are two lax monoidal functors.
  A \emph{(lax) monoidal natural transformation} is a natural transformation \(\varphi \colon F \nt G\)
  such that the following two diagrams commute for all \(x,y \in \cc\):
  \[
    \begin{mytikzcd}[ampersand replacement=\&]
      {F x \otimes F y} \& {F(x \otimes y)} \&\& 1 \& {F 1} \\
      {Gx \otimes  Gy} \& {G(x \otimes y)} \&\& G1
      \arrow["{F_{2;x,y}}", from=1-1, to=1-2]
      \arrow["{\varphi_{x \otimes y}}", from=1-2, to=2-2]
      \arrow["{\varphi_x \otimes \varphi_y}"', from=1-1, to=2-1]
      \arrow["{G_{2;x,y}}"', from=2-1, to=2-2]
      \arrow["{G_0}"', from=1-4, to=2-4]
      \arrow["{F_0}", from=1-4, to=1-5]
      \arrow["{\varphi_1}", from=1-5, to=2-4]
    \end{mytikzcd}
  \]

  Dually, given two oplax monoidal functors \(F, G \colon \cat{C} \to \cat{D}\),
  one defines an \emph{oplax monoidal natural transformation} as an arrow \(\varphi \colon F \nt G\)
  making the following diagrams commute for all \(x, y \in \cat{C}\):
  \[
    \begin{mytikzcd}[ampersand replacement=\&]
      {F(x \otimes y)} \& {F x \otimes F y} \&\& {F 1} \& 1 \\
      {G(x \otimes y)} \& {Gx \otimes  Gy} \&\& G1
      \arrow["{F_{2;x,y}}", from=1-1, to=1-2]
      \arrow["{\varphi_{x \otimes y}}"', from=1-1, to=2-1]
      \arrow["{\varphi_x \otimes \varphi_y}", from=1-2, to=2-2]
      \arrow["{G_{2;x,y}}"', from=2-1, to=2-2]
      \arrow["{G_0}"', from=2-4, to=1-5]
      \arrow["{F_0}", from=1-4, to=1-5]
      \arrow["{\varphi_1}"', from=1-4, to=2-4]
    \end{mytikzcd}
  \]
\end{definition}

\begin{definition}\label{def:lax-monoidal-adjunction}
  An adjunction \(\stdadj\) is called \emph{lax monoidal} if \(F\) and \(U\) are lax monoidal functors,
  and the unit and counit of the adjunction are monoidal natural transformations.
\end{definition}

Analogously to \cref{def:lax-monoidal-adjunction} one defines \emph{oplax monoidal adjunctions}.

\begin{definition}\label{def:lift-to-monoidal-adjunction}
  Let \(\stdadj\) be an ordinary adjunction between monoidal categories.
  A \emph{lift of \(F \dashv U\) to a lax monoidal adjunction} consists of choices of lax monoidal structures on \(F\) and \(U\),
  together with assertions that the unit and counit of the adjunction are lax monoidal natural transformations.
  In other words, it is the structure necessary for a lax monoidal adjunction
  whose left adjoint is \(F\) and whose right adjoint is \(U\).

  Similarly, one can define a lift of \(F \dashv U\) to an oplax monoidal adjunction.
\end{definition}

\begin{remark}
  A strong monoidal functor is called a \emph{monoidal equivalence} if it additionally is an equivalence of categories.
  The quasi-inverse \(F^{-1}\from \cat{D}\to \cat{C}\) of a monoidal equivalence \(F\from \cat{C} \to \cat{D}\)
  is again monoidal, and there are monoidal natural isomorphisms \(F \circ F^{-1} \iso \Id_{\cat{D}}\) and \(F^{-1} \circ F \iso \Id_{\cat{C}}\).
\end{remark}

Unless otherwise specified,
all monoidal categories in the rest of this thesis are assumed to be strict.
This is justified by the following \emph{strictification} result of Mac Lane,~\cite{MacLane63naturalassociativity},
see also~\cite[Theorem~2.8.5]{Etingof2015}.

\begin{theorem}\label{thm:monoidal-strictness-theorem}
  Every monoidal category is monoidally equivalent to a strict monoidal category.
\end{theorem}

\begin{definition}\label{def:right-module-category}
  A \emph{(left) \(\cat{C}\)\hyp{}module category}
  of a monoidal category \(\cat{C}\) consists of a category \(\cat{M}\)
  together with a functor \(\lact \from \cat{C} \times \cat{M} \to \cat{M}\)
  and isomorphisms
  \[
    \cat{M}_{\mathsf{a}; x, y, m} \from (x \otimes y) \lact m \Iso x \lact (y \lact m),
  \]
  natural in \(m \in \cat{M}\) and \(x,y \in \cat{C}\),
  satisfying \emph{coherence} axioms similar to those of a monoidal category;
  see~\cite[Definition~7.1.1]{Etingof2015}.
\end{definition}

Analogously,
one can define \emph{right} module categories over \(\cat{C}\),
which involve a functor \(\ract \from \cat{M} \times \cat{C} \to \cat{M}\) and isomorphisms
\(m \ract (x \otimes y) \Iso (x \ract y) \ract m\).
For brevity, we may sometimes speak of just ``\(\cat{C}\)\hyp{}module categories'';
this will always mean ``left \(\cat{C}\)\hyp{}module categories''.

\begin{example}[{\cite{Hajac2004a}}]\label{ex:ayd-def-as-modules-over-yd}
  Let \(H\) be a Hopf algebra with invertible antipode \(S\).
  Recall from \cref{ex:yd-modules-def-and-monoidal} that the category \(\YDM{H}\) of Yetter–Drinfeld modules is a monoidal category.
  A left-left \emph{anti\hyp{}Yetter–Drinfeld module} is some \(M \in \kVect\)
  equipped with a left \(H\)\hyp{}module structure \(\lact\)
  and a left \(H\)\hyp{}comodule structure \(m \mapsto m_{(-1)} \otimes m_{(0)}\)
  such that
  the following holds for all \(h \in H\) and \(m \in M\):\marginnote{%
    Notice that the only difference to \cref{eq:yd-condition} is that we use \(S^{-1}\) instead of \(S\).%
  }%
  \[
    {(h \lact m)}_{(-1)} \otimes {(h \lact m)}_{(0)} = h_{(1)}m_{(-1)}S^{-1}(h_{(3)}) \otimes (h_{(2)} \lact m_{(0)})
  \]

  The category  of anti\hyp{}Yetter–Drinfeld modules do not form a monoidal category, but rather a \(\cat{C}\)\hyp{}module category over the Yetter–Drinfeld modules.
  Given \(M \in \aYDM{H}\) and \(Y \in \YDM{H}\),
  define \(Y \lact M \in \aYDM{H}\) as the vector space \(Y \kotimes M\),
  equipped with the following action and coaction:
  \[
    h \lact (y \otimes m) \defeq (h_{(1)} \lact y) \otimes (h_{(2)} \lact m),
    \qquad
    y \otimes m \mapsto y_{(-1)}m_{(-1)} \otimes y_{(0)} \otimes m_{(0)}.
  \]
\end{example}

\begin{definition}\label{def:bimodule-category}
  Let \(\cat{C}\) and \(\cat{D}\) be monoidal categories.
  A \emph{\((\cat{C}, \cat{D})\)\hyp{}bimodule category} consists of a left \(\cat{C}\)\hyp{}module category \(\cat{M}\) that
  is simultaneously a right \(\cat{D}\)\hyp{}module category, such that there exists a natural isomorphism
  \[
    (x \lact m) \ract y \Iso x \lact (m \ract y),
  \]
  for all \(x \in \cat{C}\), \(y \in \cat{D}\), and \(m \in \cat{M}\),
  called the \emph{middle interchange},
  which satisfies appropriate \emph{associativity} constraints;
  see~\cite[Definition~7.1.7]{Etingof2015}.
\end{definition}

\begin{example}\label{ex:monoidal-is-bimodule}
  Every monoidal category \(\cat{C}\) is a \((\cat{C},\cat{C})\)\hyp{}bimodule category,
  whose actions are both given by the tensor product of \(\cat{C}\).
  We call this module category the \emph{regular bimodule}, and denote it by \(\leftidx{_{\cat{C}}}{\cat{C}\,}{_{\cat{C}}}\).
  Analogously, one defines the left and right regular \(\cat{C}\)\hyp{}module categories \(\leftidx{_{\cat{C}}}{\cat{C}}{}\)
  and \(\leftidx{}{\cat{C}\,}{_{\cat{C}}}\).
\end{example}

\begin{example}\label{ex:module-twisting}
  Let \(\cat{C}\) be a monoidal category.
  The regular action is not the only way in which we can consider \(\cat{C}\) as a bimodule over itself.
  Suppose that \(R\from \cat{C} \to \cat{C}\) is a strong monoidal functor.
  The action
  \[
    \ract \from \cat{C} \times \cat{C} \xrightarrow{\cat{C} \times R} \cat{C} \times \cat{C} \xrightarrow{\ \otimes\ } \cat{C}
  \]
  endows \(\cat{C}\) with the structure of a right module category over itself,
  which we shall denote by \(\cat{C}_{R}\).
  In other words, we have
  \[
    x \lact y \defeq Rx \otimes y, \qquad\qquad f \lact g \defeq Rf \otimes g,
  \]
  for \(v, w, x, y \in \cat{C}\) and \(f \from v\to w, g \from x \to y\).

  We can also twist the action from the left with another monoidal functor \(L\).
  The \emph{twisted} bimodule obtained in this manner is denoted by \({}_{L\!}{\cat{C}}_{R}\).
\end{example}

\begin{remark}\label{rem: other_variants_of_twisted_centres}
  One can easily imagine a more involved setting than \cref{ex:module-twisting}
  by twisting with an oplax monoidal functor
  \(L \from \cat{C} \to \cat{C}\) from the left
  and a lax monoidal functor \(R\) from the right.
  In this setting, \({}_{L\!}{\cat{C}}_{R}\) is a lax left and oplax right \(\cat{C}\)\hyp{}bimodule category,
  see for example~\cite[Section 2]{Szlachanyi2012:skew-monoidal}.
  Since in most of this thesis we only consider the strong case,
  we refrain from more formally treating this case;
  however, see \cref{rem: twisted_centre_of_Hopf_monad_is_rigid}.
\end{remark}

\begin{definition}\label{def:lax-module-functor}
  Let \(\cat{M}\) and \(\cat{N}\) be left module categories over a monoidal category \(\cat{C}\).
  A \emph{lax \(\cat{C}\)\hyp{}module functor} from \(\cat{M}\) to \(\cat{N}\) is a functor
  \(F \from \cat{M} \to \cat{N}\), together with a collection of morphisms
  \[
    F_{\mathsf{a};x,m} \from x \lact Fm \to F(x \lact m), \qquad \text{natural for all \(x \in \cat{C}\) and \(m \in \cat{M}\)},
  \]
  satisfying \emph{associativity} and \emph{unitality} conditions, see~\cite[Definition~7.2.1]{Etingof2015}.

  \emph{Oplax} and \emph{strong} \(\cat{C}\)\hyp{}module functors are defined similarly;
  in the former case,
  one considers arrows \(F_{\mathsf{a}; x, m} \from F(x \lact m) \to x \lact Fm\),
  while for the latter \(F_{\mathsf{a}}\) should be invertible.
\end{definition}

\begin{definition}\label{def:module-transformation}
  Let \(F, G \from \cat{M} \to \cat{N}\) be \(\cat{C}\)\hyp{}module functors between the left \(\cat{C}\)\hyp{}module categories \(\cat{M}\) and \(\cat{N}\).
  A natural transformation \(\phi \from F \nt G\) is called a
  \emph{\(\cat{C}\)\hyp{}module transformation}
  if
  \[
    G_{\mathsf{a};x, m} \circ (\phi_x \lact m) = \phi_{x \lact m} \circ F_{\mathsf{a};x, m}, \qquad\qquad
    \text{for all \(x \in \cat{C}\) and \(m \in \cat{M}\).}
  \]
\end{definition}

\begin{notation}
  For \(\cat{C}\)\hyp{}module categories \(\cat{M}\) and \(\cat{N}\),
  we obtain the following categories of \(\cat{C}\)\hyp{}module functors:
  \begin{itemize}
    \item \(\cat{C}\mathsf{Mod}(\cat{M},\cat{N}) \!\defeq\! \mathsf{Str}\cat{C}\mathsf{Mod}(\cat{M,N})\):
    strong \(\cat{C}\)\hyp{}module functors from \(\cat{M}\) to \(\cat{N}\),
    \item \(\mathsf{Lax}\cat{C}\mathsf{Mod}(\cat{M,N})\):
    lax \(\cat{C}\)\hyp{}module functors from \(\cat{M}\) to \(\cat{N}\),
    \item \(\mathsf{Oplax}\cat{C}\mathsf{Mod}(\cat{M,N})\):
    oplax \(\cat{C}\)\hyp{}module functors from \(\cat{M}\) to \(\cat{N}\).
  \end{itemize}
\end{notation}

\begin{example}\label{ex:regularend}
  For an object \(x\) is a (non-strict) monoidal category \(\cat{C}\),
  the functor \(F \defeq \blank \otimes x\) is a strong left \(\cat{C}\)\hyp{}module functor.
  The transformation \(F_{\mathsf{a}}\) is given by the associator of \(\cat{C}\).
  Further, a morphism \(f\from x \to y\) in \(\cat{C}\) yields a \(\cat{C}\)\hyp{}module transformation \(\blank \otimes f\from \blank \otimes x \nt \blank \otimes y\).
\end{example}

In fact, \cref{ex:regularend} generalises and completely describes module functors from the regular module.
We emphasise that this result is a consequence of the bicategorical Yoneda lemma,
see~\cite[Lemma~8.3.16]{johnson-yau2021:TwoDimensionalCategories}.

\begin{proposition}\label{prop:bicategorical-yoneda-correspondence-monads-and-algebras}
  Let \(\cat{M}\) be a left module category over the monoidal category \(\cat{C}\).
  Then there is an equivalence of module categories
  \[
    \mathsf{Str}\cat{C}\mathsf{Mod}(\cat{C}, \cat{M}) \iso \cat{M}, \qquad \qquad
    F \mapsto F1, \qquad \qquad
    \blank \lact m \longmapsfrom m.
  \]
\end{proposition}

In particular, this yields a monoidal equivalence \(\mathsf{Str}\cat{C}\mathsf{Mod} \simeq \cat{C}^{\tensorop}\).

\subsection{Braidings}\label{subsec: braided_and_ribbon_cats}

\textsc{Braidings are natural transformations} relating the tensor product to its opposite.
They where introduced by Joyal and Street in~\cite{joyal85:braided,joyal86:braided,Joyal1993},
and build on the notion of symmetries studied in~\cite{MacLane63naturalassociativity,Eilenberg1966}.

\begin{definition}\label{def: braided_cat}
  A \emph{braiding} on a monoidal category \(\cat{C}\) is a natural isomorphism
  \(\sigma_{\blank, \bblank} \from \blank \otimes \bblank \nt \bblank \otimes \blank\),
  satisfying \(\sigma_{x,1} = \id_x\) and the \emph{hexagon axioms}\note{
    The name ``hexagon axioms'' is due to the fact
    that in the non-strict setting,
    the defining equations
    can be organised as a hexagon-shaped diagram, see~\cite{Joyal1993}.
  }.

  The pair \((\cat{C}, \sigma)\) will be called a \emph{braided monoidal category}.
\end{definition}

\begin{remark}\label{rem: braiding_compatibility_with unit}
  For a braiding \(\sigma\), the condition that \(\sigma_{x,1} = \id_x\) is actually implied by the hexagon identities:
  The calculation
  \[
    \sigma_{x,1} = \sigma_{x, 1\otimes 1} = (\id_1 \otimes \sigma_{x,1}) \circ (\sigma_{x,1} \otimes \id_1) = \sigma_{x,1} \circ \sigma_{x,1}
  \]
  shows that \(\sigma_{x,1}\) is an invertible idempotent,
  hence the claim follows.
\end{remark}

\begin{example}[{\cite[Theorem~7.2]{yetter90:quant}}]\label{ex:yd-modules-are-braided}
  The category of left-left Yetter–Drinfeld modules of \cref{ex:yd-modules-def-and-monoidal}
  is braided monoidal when equipped with
  \[
    \sigma \defeq {\{\, \sigma_{Y,Z} \from Y \otimes Z \to Z \otimes Y, \quad y \otimes z \mapsto (y_{(-1)} \lact z) \otimes y_{(0)} \,\}}_{Y,Z \,\in\, \YDM{H}},
  \]
  called the \emph{\YetterDrinfeld{} braiding}.
  Note that \(\sigma\) is invertible because
  \(H\) has an invertible antipode.
  Its inverse is given for all \(Y, Z \in \YDM{H}\) by
  \[
    \sigma^{-1} \defeq {\{\, \sigma^{-1}_{Y,Z} \from Z \otimes Y \to Y \otimes Y, \quad z \otimes y \mapsto y_{(0)} \otimes (S^{-1}(y_{(-1)}) \lact z) \,\}}_{Y,Z \,\in\, \YDM{H}}.
  \]
\end{example}

\begin{definition}\label{def:symmetric-braiding}
  A braiding \(\sigma\) on a monoidal category \(\cat{C}\) is called \emph{symmetric} if \(\sigma_{x,y}^{-1} = \sigma_{y,x}\) for all \(x, y \in \cat{C}\).
  A monoidal category equipped with a symmetric braiding will be referred to as a \emph{symmetric monoidal category}.
\end{definition}

Braidings are depicted in the graphical calculus by crossings of strings subject to Reidemeister\hyp{}esque identities, see~\cite{Selinger2011}.
The following figure shows a braiding, its inverse, the hexagon identity, and the naturality of the braiding in its first argument:
\[
  \tikzfig{braiding-properties}
\]

\subsection{Closedness}\label{sec:closed-mon-cats}

\begin{definition}\label{def:closed-module-category}
  Let \(\cat{C}\) be a monoidal category.
  An object \(m\) in a \(\cat{C}\)\hyp{}module category \(\cat{M}\) is called \emph{closed}
  if the functor \(\blank \triangleright m\from \cat{C} \to \cat{M}\) has a right adjoint
  \(\hom{m, \blank}\from \cat{M}\to \cat{C}\), called the \emph{internal hom from \(m\)}.
  We refer to
  \[
    \adj{\blank \lact m}{\hom{m, \blank}}{\cc}{\mm}
  \]
  as the \emph{internal tensor--hom adjunction}.
\end{definition}

Similarly to \cref{def:closed-module-category},
an object \(m \in \cat{M}\) is called \emph{coclosed} if the functor \(\blank \triangleright m\from \cat{C} \to \cat{M}\) has a left adjoint;
we denote this \emph{internal cohom from \(m\)} by \(\cohom{m, \blank}\from \cat{M}\to \cat{C}\),
and call \(\adj{\cohom{m, \blank}}{\blank \lact m}{\mm}{\cc}\) as the \emph{internal cohom--tensor adjunction}.
If every object of \(\cat{M}\) is closed,
we call \(\cat{M}\) a \emph{closed \(\cat{C}\)\hyp{}module category}.

\begin{definition}\label{def:left-closed-monoidal-category}
  A monoidal category \(\cat{C}\) is called \emph{left closed}
  if the left regular \(\cat{C}\)\hyp{}module category \({}_{\cat{C}}\cat{C}\) is closed,
  \emph{right closed} if the right regular \(\cat{C}\)\hyp{}module category \(\cat{C}_{\cat{C}}\) is closed,
  and \emph{closed} or \emph{biclosed} if it is both left and right closed.

  In these cases, we write \({[x, \blank]}_{\ell}\) and \({[x, \blank]}_r\) for \(\hom{x, \blank}\), respectively.
\end{definition}

\begin{example}
  A symmetric monoidal category is left closed if and only if it is right closed if and only if it is closed.
\end{example}

The object-wise tensor--hom adjunctions of a closed \(\cat{C}\)\hyp{}module category specify a unique functor
\(\hom{\blank, \bblank} \from \cat{M}^{\op}\times\cat{M} \to\cat{C}\) such that we have
\[
  \cat{M}(c \lact m, n) \cong \cat{C}(c, \hom{m, n}),
\]
natural in all three variables, see~\cite[Section~\textsc{iv}.7]{MacLane1998} for the monoidal case.

\begin{example}\label{ex:deligne-morphism-in-closed-monoidal-cat}
  For any object \(x \in \cat{C}\) of a right closed monoidal category,
  a combination of the unit and counit
  \[
    \eta^{(x)}_y \colon y \to {[x, x \otimes y]}_r, \qquad \qquad
    \varepsilon^{(x)}_y \colon x \otimes {[x,y]}_r \to y, \qquad \text{for all } y \in\cat{C}
  \]
  gives rise to the morphism
  \begin{equation}\label{eq:canonical-morphism-deligne}
    \begin{mytikzcd}[ampersand replacement=\&]
      {\phi^{(x)}_y \defeq {[x,1]}_r\otimes y} \& {{[x, x \otimes {[x,1]}_r\otimes y]}_r} \& {{[x,y]}_r.}
      \arrow["{\eta^{(x)}_{{[x,1]}_r \otimes y}}", from=1-1, to=1-2]
      \arrow["{[x, \varepsilon^{(x)}_1 \otimes y]}", from=1-2, to=1-3]
    \end{mytikzcd}
  \end{equation}
\end{example}

\subsection{Rigidity and pivotality}\label{subsubsec: rigid_and_piv_cats}

\textsc{Rigidity}\marginnote{%
  Rigid monoidal categories are also called \emph{autonomous},
  or, in the symmetric case, \emph{compact closed}.%
}
\textsc{in the context of monoidal categories} refers to a concept of duality similar to that of finite\hyp{}dimensional vector spaces.
Importantly, notions like dual bases and evaluations have their analogues in this setting.

Pivotal categories%
—also known as \emph{balanced} or \emph{sovereign} categories—%
are those where there exists an identification between objects and their double duals that is compatible with the tensor product.

Recall from \cref{eq:canonical-morphism-deligne}
there is a canonical morphism
\[
  \Phi^{(x)}_y \from {[x, 1]}_r \otimes y \to {[x, y]}_r
\]
for all objects \(x, y\) in the right closed monoidal category \(\cat{C}\).
The next results links the invertibility of this morphism to conditions that \(\cat{C}\) is rigid monoidal;
see for example~\cite{kelly72:many}, and~\cite[Proposition~2.1]{niefield17:coexp} for a concise proof.

\begin{proposition}\label{prop:dualisable-conditions-Deligne}
  Let \(\cat{C}\) be a monoidal category.
  For any object \(x \in \cat{C}\),
  the following statements are equivalent:
  \begin{thmlist}[itemsep=1ex]
    \item the right internal hom of \(x\) exists, and
    the canonical arrows \(\phi^{(x)}_{y}\) are invertible for all \(y \in\cat{C}\);
    \item the right internal hom of \(x\) exists and \(\phi^{(x)}_x\from {[x, 1]}_r \otimes x \to {[x, x]}_r\) is an isomorphism; and
    \item there exists an object \(\rd{x} \in\cat{C}\) together with morphisms \(\ev^r_x\colon x \otimes \rd{x} \to 1\) and \(\coev^r_x\colon 1 \to \rd{x} \otimes x\), satisfying the snake identities
    \begin{equation}\label{eq:snake-ids}
      \begin{aligned}
        \id_x &= (\ev^r_x \otimes \id_x) \circ (\id_x \otimes \coev^r_x), \\
        \id_{\rd{x}} &= (\id_{\rd{x}} \otimes \ev^r_x) \circ (\coev^r_x \otimes \id_{\rd{x}}).
      \end{aligned}
    \end{equation}
  \end{thmlist}
\end{proposition}

\begin{definition}\label{def:rigidly-dualisable}
  We call \(x\) \emph{right (rigidly) dualisable}
  if any of the equivalent conditions of \cref{prop:dualisable-conditions-Deligne} are met.
  In this case, we have \({[x, y]}_r \cong \rd{x} \otimes y\) for all \(x, y \in \cat{C}\) and in particular \(\rd{x} \cong {[x, 1]}_r\).
  The \emph{left (rigid) dualisability} of an object \(x\in\cat{C}\) can be defined similarly by either the invertibility of a canonical morphism
  \(\psi^{(x)}_y \colon y \otimes {[x,1]}_{\ell} \to {[x,y]}_{\ell}\) or by the existence of an object \(\ld{\!{}x} \in\cat{C}\)
  endowed with morphisms \(\ev^{\ell}_x \colon \ld{\!{}x} \otimes x \to 1\) and \(\coev^{\ell}_x \colon 1 \to x \otimes \ld{\!{}x}\),
  subject to suitable variants of the snake identities.
\end{definition}

When the context leaves the choice unambiguous,
we—in the interest of brevity—omit the left and right superscripts for the (co)evaluation morphisms.
The pair \((y, x, \ev, \coev)\),
where \(y\) is a left dual of \(x\) and, thus, \(x\) is a right dual of \(y\)
is called a \emph{dual pair}.

\begin{definition}
  A monoidal category \(\cat{C}\) is called \emph{left (right) rigid}
  if every object is left (right) rigidly dualisable,
  and \emph{rigid} if it is left and right rigid.

  We denote the left and right \emph{dualising functors} of a rigid category \(\cat{C}\) by
  \[
    \ld{\!(\blank)} \cong {[\blank, 1]}_{\ell}\from \cat{C}^{\op,\tensorop} \to \cat{C}
    \qquad\text{and}\qquad
    \rd{(\blank)} \cong {[\blank, 1]}_{r}\from \cat{C}^{\op,\tensorop} \to \cat{C}.
  \]
\end{definition}

\begin{example}\label{ex:left-dualising-functor}
  Let \(\cat{C}\) be a left rigid monoidal category.
  Given a morphism \(f \from x \to y\) in \(\cat{C}\),
  the induced morphism \(\ld{\!{}f}\from \ld{\!{}y} \to \ld{\!{}x}\) is given by
  \[
    (\ev_y^{\ell} \otimes \id_{\ld{\!{}x}}) \circ (\id_{\ld{\!{}y}} \otimes f \otimes \id_{\ld{\!{}x}}) \circ (\id_{\ld{\!{}y}} \otimes \coev_x^{\ell}).
  \]
\end{example}

There are categories with only one-sided closedness or rigidity, see~\cite[Theorem~6.3.3]{loregian2021} and~\cite[Example~1.6.2]{Turaev2017}.
Taking the opposite of a rigid category reverses the roles of evaluation and coevaluation.
Thus the opposite \(\cat{C}^{\op}\) of the left rigid category \(\cat{C}\) is right rigid.

\begin{lemma}[{\cite[Exercise~2.10.6]{Etingof2015}}]\label{lem:transfer-of-rigidity}
  Let \(F \from \cat{C} \to \cat{D}\) be a strong monoidal functor. The image \(Fx\) of  any (rigidly) dualisable object \(x \in\cat{C}\) is dualisable.
\end{lemma}
\begin{proof}[Sketch of proof]
  If \(x \in \cat{C}\) is an object with right dual \(\rd{x}\),
  define \(\rd{(Fx)} \defeq F(\rd{x})\);
  the evaluation map is given by
  \[
    Fx \otimes F(\rd{x}) \xrightarrow{\ F_2\ } F(x \otimes \rd{x}) \xrightarrow{F \ev^{(r)}_x} F1 \xrightarrow{F_0^{-1}} 1,
  \]
  and the coevaluation is defined analogously.
  A straightforward calculation shows the snake identities to be satisfied.
\end{proof}

Graphically, evaluations and coevaluations will be represented by semicircles,
possibly decorated with arrows to emphasise whether we consider their left or right version.
\[
  \tikzfig{rigid-in-string-diagrams}
\]

\begin{definition}\label{def: invertible_objects}
  An object \(x \in \cat{C}\) in a rigid monoidal category \(\cat{C}\) is called \emph{invertible} if its (left) evaluation and coevaluation are isomorphisms.
\end{definition}

It follows that the right evaluations and coevaluations of an invertible objects are isomorphisms as well.
Further, tensor products and duals of invertible objects are invertible,
such that the full subcategory \(\mathsf{Inv}(\cat{C})\) of invertible object of \(\cat{C}\) is rigid monoidal.

\begin{definition}[{\cite{castelnuovo1905}}]\label{def: invertible_object}
  The \emph{Picard group} \(\Pic \cat{C}\) of a rigid monoidal category \(\cat{C}\) is the group of isomorphism classes of invertible objects in \(\cat{C}\).
  Its multiplication is induced by the tensor product; \ie, \([\alpha] \cdot [\beta] \defeq [\alpha \otimes \beta]\) for \(\alpha, \beta \in \mathsf{Inv}(\cat{C})\).
  The unit of \(\Pic \cat{C}\) is \([1]\) and for any \(\alpha \in \mathsf{Inv}(\cat{C})\) we have that \({[\alpha]}^{-1} =[\ld{\!\alpha}]\).
\end{definition}

\begin{proposition}\label{prop: rigidity_yields_adjoints}
  For every object \(x \in \cat{C}\) in a rigid category \(\cat{C}\) we obtain two chains of adjoint endofunctors of \(\cat{C}\):
  \begin{align*}
    \dots \adjoint \blank \otimes \rd{x} \adjoint \blank \otimes x \adjoint \blank \otimes \ld{\!{}x} \adjoint \dots
    \qquad
    \dots \adjoint \ld{\!{}x} \otimes \blank \adjoint x \otimes \blank \adjoint \rd{x} \otimes \blank \adjoint \dots
  \end{align*}
  Further, \(\blank \otimes x\) and \(x \otimes \blank\) are equivalences of categories if and only if \(x\) is invertible.
\end{proposition}
\begin{proof}
  The existence of the stated chains of adjunctions follows from~\cite[Proposition~2.10.8]{Etingof2015}.
  A straightforward calculation shows that tensoring (from the left or right)
  with an invertible object yields an equivalence of categories.

  Conversely, suppose that \(x \in \cat{C}\) is such that \(F \defeq \blank \otimes x\) is an equivalence of categories.
  The functor \(F\) and its quasi-inverse \(U\) are part of an adjunction
  with invertible unit \(\eta \from {\Id_{\cat{C}}} \nt UF\)
  and counit \(\varepsilon \from FU \nt \Id_{\cat{D}}\),
  see for example~\cite[Proposition~4.4.5]{category-theory-in-context_riehl}.
  By~\cite[Proposition~4.4.1]{category-theory-in-context_riehl}, there exists a natural isomorphism \(\theta \from U \nt \blank \otimes \ld{\!{}x}\)
  that commutes with the respective counits and units.
  Applied to the monoidal unit \(1 \in \cat{C}\), we obtain
  \[
    \coev_{x}^{\ell} = \theta_x \circ \eta_1 \qquad \text{and} \qquad \ev_x^{\ell} \circ (\theta_1 \otimes \id_x) = \varepsilon_1 .
  \]
  It follows that \(x\) is invertible.
  An analogous argument shows that \(x \otimes \blank \) being an equivalence of categories also entails \(x\) being invertible.
\end{proof}

\begin{remark}\label{rmk:trep-is-not-rigid-first-exposure}
  Given a monoidal category \(\cat{C}\),
  the mere presence of adjunctions
  \[
    \adj{\blank \otimes x}{\blank \otimes Lx}{\cat{C}}{\cat{C}}
    \qquad\qquad\text{and}\qquad\qquad
    \adj{x \otimes \blank}{ Rx \otimes \blank}{\cat{C}}{\cat{C}}
  \]
  for objects \(Lx, Rx \in \cat{C}\) does not lead to rigidity,
  but to the weaker notion of \emph{tensor representability}, see \cref{def:tensor-representable-endo-functors,def:tensor-representable}.

  To elucidate the underlying problem,
  let us assume for a moment that we are given objects \(x, Rx \in \cat{C}\),
  such that \(\adj{x \otimes \blank }{Rx \otimes\blank}{\cat{C}}{\cat{C}}\).
  One can show that any right rigid dual of \(x\)—if it exists—has to be isomorphic to \(Rx\);
  \ie, one has \(Rx \cong \rd{x}\) on objects.
  Hence, the unit \(\eta^{(x)}_{z} \from z \to Rx \otimes x \otimes z\) and counit \(\varepsilon^{(x)}_{z} \from x \otimes Rx \otimes z \to z \) of the adjunction provide us with natural candidates for the coevaluation and evaluation morphisms:
  \[
    \coev_x \defeq \eta^{(x)}_1 \from 1 \to Rx \otimes x \qquad \qquad \text{and} \qquad \qquad \ev_x \defeq \varepsilon^{(x)}_1 \from x \otimes Rx \to 1.
  \]
  Evaluating the triangle identities of the adjunction at the monoidal unit yields
  \[
    \id_{x} = \varepsilon^{(x)}_{x} \circ (x \otimes \eta^{(x)}_1)
    \qquad\qquad\text{and}\qquad\qquad
    \id_{Rx} = (Rx \otimes \varepsilon^{(x)}_1) \circ \eta^{(x)}_{Rx}.
  \]
  However, if \(Rx\) is to be a dual of x in the rigid sense,
  the snake identities~\eqref{eq:snake-ids} must hold.
  For this, we ought to require the stronger condition that
  \[
    \varepsilon^{(x)}_x = \varepsilon^{(x)}_1 \otimes \id_x
    \qquad\qquad\text{and}\qquad\qquad
    \eta^{(x)}_{Rx} = \eta^{(x)}_1 \otimes Rx.
  \]
\end{remark}

\begin{remark}
  A left closed monoidal such that for every \(x \in \cat{C}\),
  the functor \(\blank \otimes x\)
  admits a right dual \({[x, \blank]}_{\ell}\) as a \(\cat{C}\)\hyp{}module functor%
  —see \cref{ex:module-adjunction-as-comodule-adjunction}—%
  is already rigid.
  Explicitly, this involves the following map natural in \(y, z \in \cat{C}\):
  \[
    {[x, z]}_{\ell} \otimes y \iso {[x, z \otimes y]}_{\ell},
  \]
  the resulting compatibility conditions of which imply the snake identities.
\end{remark}

We will study these questions in greater details in \cref{sec:trep-categories};
see also \cref{rmk:rigid-iff-lax-is-strong}
for a characterisation of rigidity in terms of \(\cat{C}\)\hyp{}module functors.

\begin{definition}\label{def: stric_rigid_cat}
  A rigid monoidal category \(\cat{C}\) is called \emph{strict rigid}
  if the dual functors \(\ld{\!(\blank)},\,\rd{(\blank)}\from \cat{C}^{\op,\tensorop} \to \cat{C}\) are strict and
  \[
    \ld{\!\big(\rd{(\blank)}\big)} = \Id_{\cat{C}}
    = \rd{\big(\ld{\!(\blank)} \big)\!}.
  \]
\end{definition}

\begin{notation}\label{almostlisp}
  Let \(\cat{C}\) be a rigid monoidal category;
  for \(x \in \cat{C}\) and \(n \in \integer\), write
  \[
    {(x)}^{n} \defeq
    \begin{cases}
      \text{The \(n\)\hyp{}fold left dual of \(x\)}, & \text{if } n > 0; \\
      x, & \text{if } n = 0; \\
      \text{The \(n\)\hyp{}fold right dual of \(x\)}, & \text{if } n < 0.
    \end{cases}
  \]
\end{notation}

Our next result was conjectured in~\cite[Section 5]{Schauenburg2001},
and is a slight variation of~\cite[Theorem~2.2]{Ng2007}.
It shows that every rigid category admits a \emph{rigid strictification};
\ie, a monoidally equivalent strict rigid category.
The compatibility between the respective left and right duality functors is an immediate consequence of the fact that for any strong monoidal functor \(F \from \cat{C} \to \cat{D}\) between rigid categories there are natural monoidal isomorphisms
\[
  \phi_x \from F (\ld{\!x}) \iso \ld{Fx},
  \qquad\text{and}\qquad
  \psi_x \from F (\rd{x}) \iso \rd{Fx},
  \quad \text{ for all } x \in \cat{C}.
\]

\begin{theorem}\label{thm: rigid_strictification}
  Every rigid category admits a rigid strictification.
\end{theorem}
\begin{proof}
  Suppose that \(\cat{C}\) is a rigid strict monoidal category \(\cat{C}\).
  Build a monoidally equivalent strict rigid category \(\cat{D}\) as follows:
  the objects of \(\cat{D}\) are (possibly empty) finite sequences
  \((x_1^{n_1}, \dots, x_i^{n_i})\) of objects \(x_1, \dots, x_i \in \cat{C}\),
  adorned with integers \(n_1, \dots, n_i \in \integer\).

  To define the morphisms of \(\cat{D}\),
  recall \cref{almostlisp} and set
  \[
    \cat{D}(
    (x_1^{n_1}, \dots, x_i^{n_i}),
    (y_1^{m_1}, \dots, y_j^{m_j})
    )
    \defeq
    \cat{C}(
    {(x_1)}^{n_1} \otimes \dots \otimes {(x_i)}^{n_i},
    {(y_1)}^{m_1} \otimes \dots \otimes {(y_j)}^{m_j}
    ).
  \]
  The category \(\cat{D}\) is strict monoidal when equipped with the concatenation of sequences as tensor product and the empty sequence as unit.
  By construction, there exists a strict monoidal equivalence of categories \(F \from \cat{D} \to \cat{C}\),
  which maps any object
  \((x_1^{n_1}, \dots, x_i^{n_i}) \in \cat{D}\) to \({(x_1)}^{n_1} \otimes \dots \otimes {(x_i)}^{n_i} \in \cat{C}\),
  as well as every morphism to itself;
  the unit of \(\cat{C}\) is the empty tensor product.

  Define the left dual of an object
  \(x \defeq (x_1^{n_1}, \dots, x_i^{n_i}) \in \cat{D}\)
  as
  \[
    \ld{\!{}x} \defeq (x_i^{n_i+1}, \dots, x_1^{n_1+1}),
  \]
  with
  evaluation and coevaluation morphisms given by
  \[
    \tikzfig{rigid-strict-ev-coev}
  \]
  For all \(1 \leq k \leq i\), we set
  \[
    \phi_k \defeq
    \begin{cases}
      \ev^{\ell}_{{(x_k)}^{n_k}}, &\text{ if } n_k \geq 0; \\
      \ev^r_{{(x_k)}^{n_k+1}}, &\text{ if } n_k < 0; \\
    \end{cases}
    \qquad
    \text{ and }
    \qquad
    \psi_k \defeq
    \begin{cases}
      \coev^{\ell}_{{(x_k)}^{n_k}}, & \text{ if } n_k \geq 0; \\
      \coev^r_{{(x_k)}^{n_k+1}}, &\text{ if } n_k < 0. \\
    \end{cases}
  \]
  The right dual of \(x\) if defined similarly as \(\rd{x} \defeq (x_i^{n_i-1}, \dots, x_1^{n_1-1})\).
  Hence, \(\cat{D}\) is strict rigid, which completes the proof.
\end{proof}

Module functors over rigid categories behave in particularly nice ways.

\begin{proposition}[{\cite[Remark~4]{ostrik03:modul-hopf},~\cite[Lemma~2.10]{douglas19}}]\label{prop:lax-is-strong-when-rigid}
  If \(\cat{C}\) is a left rigid monoidal, then every lax \(\cat{C}\)\hyp{}module functor is strong.
  Dually, if \(\cat{C}\) is a right rigid monoidal category, then every oplax \(\cat{C}\)\hyp{}module functor is strong.
\end{proposition}

Many applications require that the objects of a rigid monoidal category are isomorphic to their double duals in a way which is compatible with the monoidal structure.
In \cref{cor: pivotal_from_central_anti_central} we prove a result for detecting this.

\begin{definition}\label{def: pivotal_category}
  A \emph{pivotal category} is a rigid monoidal category \(\cat{C}\) equipped with a monoidal natural isomorphism
  \(\rho\from\! \Id_{\cat{C}} \Iso\! \lbiDual{\!\blank}\),
  a \emph{pivotal structure} of \(\cat{C}\).
\end{definition}

Rigid monoidal categories do not have to admit a pivotal structure and, if they do, it need not be unique.
Examples coming from Hopf algebra theory are given in~\cite{Kauffman1993,Halbig2018,Halbig2019}.
However, every rigid monoidal category admits a universal pivotal category,
called its \emph{pivotal cover}, see~\cite{Shimizu2015}.

\subsection{The Drinfeld centre}\label{sec:drinfeld-centre}

\textsc{Classically, the centre construction} is used to build a braided monoidal category from a monoidal one,
see for example~\cite[Chapter~7]{Etingof2015}.
Throughout especially \cref{sec:twisted-centres,sec:monadic-tannaka-reconstruction},
we work in a slightly more general setting,
see for example~\cite{Gelaki2009,Bruguieres2012,Hassanzadeh2019,femić23:categ-yetter,kowalzig24:centr}.

\begin{definition}\label{def: half_braiding}
  Let \(\cat{C}\) be a monoidal category, \(\cat{M}\) a \(\cat{C}\)\hyp{}bimodule category, and \(m \in \cat{M}\).%
  \marginnote{%
    Analogously to \cref{rem: braiding_compatibility_with unit},
    one can show that
    \(\sigma_{M, 1} = \id_M\) is implied by the hexagon axiom.%
  }
  A \emph{half-braiding} on \(m\) is a natural isomorphism
  \(\sigma_{\blank, \bblank} \from \blank \ract \bblank \nt \blank \lact \bblank\),
  such that for all \(x, y \in \cat{C}\) we have \(\sigma_{M, 1} = \id_M\) and
  \[
    \sigma_{m,x\otimes y} = ({\id_x} \lact \sigma_{m,y}) \circ (\sigma_{m, x} \ract \id_y).
  \]
\end{definition}

\begin{definition}\label{def: centre_bimodue_category}
  The \emph{centre} of a \(\cat{C}\)\hyp{}bimodule category \(\cat{M}\) is the category \(\ZCat{M}\) defined as follows:
  \begin{itemize}
    \item Objects are pairs \((m, \sigma_{m, \blank})\) of an object \(m \in \cat{M}\) and a half-braiding \(\sigma_{m,\blank}\).
    \item A morphism \(f \from (m, \sigma_{m,\blank}) \to (n, \sigma_{n,\blank})\) consists of an \(f \in \cat{M}(m,n)\) that commutes with the half-braidings:
    \[
      (\id_x \lact f) \circ \sigma_{m,x} = \sigma_{n,x} \circ (f \ract \id_x), \qquad \qquad \text{ for all } x \in \cat{C}.
    \]
  \end{itemize}
\end{definition}

There is a canonical forgetful functor \(U^{(M)} \from \ZCat{M} \to \cat{M}\).
Unlike classical representation theory where the centre of a bimodule is a subset of the bimodule,
\(U^{(M)}\) need not be injective on objects in general.

\begin{example}\label{ex: Drinfeld_centre}
  The centre \(\ZCat{C}\) of the regular bimodule of a monoidal category \(\cat{C}\),
  see \cref{ex:monoidal-is-bimodule},
  is the \emph{Drinfeld centre} of \(\cat{C}\).
  The tensor product is defined by
  \((x, \sigma_{x, \blank}) \otimes (y, \sigma_{y, \blank}) \defeq ( x \otimes y, \sigma_{x\otimes y, \blank})\), with
  \[
    \sigma_{x\otimes y, z} \defeq (\sigma_{x, z} \otimes \id_y) \circ ({\id_x} \otimes \sigma_{y, z}), \qquad \qquad \text{ for all } z \in \cat{C}.
  \]
  The centre is braided monoidal, with braiding given by gluing together the respective half-braidings.
  The hexagon axioms follow from the definition of the half-braiding and the tensor product of \(\ZCat{C}\).
\end{example}

\begin{example}\label{ex:yd-modules-as-drinfeld-centre}
  Let \(H \in \kVect\) be a finite-dimensional Hopf algebra.
  Then the category of left-left Yetter–Drinfeld modules
  of \cref{ex:yd-modules-def-and-monoidal,ex:yd-modules-are-braided}
  is, as a braided monoidal category,
  equivalent to the Drinfeld centre of the category \(\Tetramod[][H]\) of left \(H\)\hyp{}modules,~\cite{drinfeld87:quant,yetter90:quant}.

  Further, by~\cite{drinfeld87:quant} there exists another Hopf algebra \(D(H)\),
  the \emph{Drinfeld double} of \(H\),
  that has \(\cent*{\Tetramod[][H]{}} \simeq \YDM{H}\) as its category of modules.
  We refer to~\cite[Definition~\textsc{ix}.4.1]{Kassel1998}
  for a complete definition,
  and to~\cite[Theorem~\textsc{xiii}.5.1]{Kassel1998}
  for a proof of the above equivalence.
\end{example}

\begin{remark}\label{rmk:anti-drinfeld-double}
  The category of anti\hyp{}Yetter–Drinfeld modules of \cref{ex:ayd-def-as-modules-over-yd}
  is also the category of modules over a certain Hopf algebra:
  the \emph{anti\hyp{}Drinfeld double}~\cite{caenepeel97:cross-doi-hopf,schauenburg99:examp-doi-koppin-hopf-yetter-drinf}.
  In analogy to how the anti\hyp{}Yetter–Drinfeld modules are a module category over the Yetter–Drinfeld modules,
  the anti\hyp{}Drinfeld double is a comodule algebra over the Drinfeld double.
\end{remark}

For the next result, recall \cref{almostlisp} for iterated duals.

\begin{proposition}[{\cite[Lemma~7]{joyal91}}]\label{prop: rigdity_of_centre}
  The centre of a (strict) rigid category \(\cat{C}\) is (strict) rigid:
  for all \((x, \sigma_{x,-}) \in \ZCat{C}\),
  we have that \(U^{(Z)} \big(\ld{(x, \sigma_{x,-})}\big) = \ld{\!{}x}\)
  and \(U^{(Z)} \big(\rd{(x, \sigma_{x,-})}\big) = \rd{x}\).
  Moreover, for every \(n \in \integer\) and \(x \in \ZCat{C}\),
  the equality \(\sigma_{{(x)}^{n}, {(y)}^{n}} = {(\sigma_{x,y})}^{n}\) holds for all \(y \in \cent{C}\).
\end{proposition}

\section{Linear and abelian categories}\label{sec:linear-ab-cats}

\textsc{We shall pay special attention}
to \(\Bbbk\)\hyp{}linear categories,
for some field \(\Bbbk\).
Indeed, \(\Bbbk\)\hyp{}linear functor categories encompass many phenomena in representation theory;
see for example \cref{sec:tensor-rep-functor-cats,sec:hopf-trimodules}.
We shall denote the category \(\Bbbk\)\hyp{}vector spaces by \(\kVect\),
and write \(\kvect\) for the full subcategory of finite\hyp{}dimensional \(\Bbbk\)\hyp{}vector spaces.\note{%
  This will be a general theme throughout the thesis:
  whenever there exists a full subcategory of ``finite\hyp{}dimensional objects''
  inside of a lager category of all objects,
  the finite\hyp{}di\-men\-sio\-nal subcategory will start with a lower case letter,
  and we use a capital for the larger category.%
}

\begin{definition}\label{def:linear-category}
  A \emph{\(\Bbbk\)\hyp{}linear category} is a category enriched in \(\kVect\)
\end{definition}

\begin{example}\label{rmk:linearisation}
  Let \(\cat{C}\) be an ordinary category.
  Its \emph{linearisation} \(\Bbbk\cat{C}\) has the same objects,
  and the space of morphisms between two objects \(a, b \in \Bbbk\cat{C}\) is \(\Bbbk \cat{C}(a, b) \defeq \mathrm{span}_{\Bbbk}\cat{C}(a,b)\).
  To define the composition in \(\Bbbk\cat{C}\),
  note that for any  \(a, b, c \in \Bbbk\cat{C}\) there is a unique linear map
  \[
    \blank \circ \bblank \from \Bbbk\cat{C}(b,c) \kotimes \Bbbk\cat{C}(a,b) \to \Bbbk\cat{C}(a,c),
  \]
  such that \(g \circ f = gf\)  for all  \(g\in \cat{C}(b,c)\) and  \(f\in \cat{C}(a,b)\).

  Let \(\iota\from \cat{C} \to \Bbbk\cat{C}\) be the functor that is the identity on objects and maps any morphism to the corresponding basis vector.
  For any functor \(F \from \cat{C} \to \cat{D}\) whose codomain is a \(\Bbbk\)\hyp{}linear category, we obtain a commuting triangle:
  \[
    \begin{mytikzcd}[ampersand replacement=\&]
      {\cat{C}} \&\& \cat{D} \\
      \& {\Bbbk\cat{C}}
      \arrow["F", from=1-1, to=1-3]
      \arrow["\iota"', from=1-1, to=2-2]
      \arrow["{\exists!\,G \; (\text{linear})}"', from=2-2, to=1-3]
    \end{mytikzcd}
  \]
  As a consequence, if we endow the ordinary functor category \([\cat{C}, \kVect]\) with the pointwise \(\Bbbk\)\hyp{}linear structure,
  we obtain an isomorphism of linear categories between it and the category of \(\Bbbk\)\hyp{}linear functors from \(\Bbbk\cat{C}\) to \(\kVect\).
\end{example}

\subsection{Finite and locally finite abelian categories}\label{sec:loc-fin-ab-categories}

\textsc{In this section},
we assume all categories and functors to be \(\Bbbk\)\hyp{}linear,
for some field \(\Bbbk\).
We call a functor \(F\from \cat{A} \to \cat{B}\) between abelian categories \(\cat{A}\) and \(\cat{B}\)
\emph{left exact} if it preserves finite limits,
\emph{right exact} if it preserves finite colimits,
and \emph{exact} if it is right and left exact.
The category of right exact functors from \(\cat{A}\) to \(\cat{B}\) will be denoted by \(\mathsf{Rex}(\cat{A}, \cat{B})\),
and we write \(\mathsf{Lex}(\cat{A}, \cat{B})\) for the category of left exact functors.
Further, let \(\proj{\cat{A}}\) denote the full subcategory of \(\cat{A}\) consisting of its projective objects,
and similarly for \(\inj{\cat{A}}\).

We recall basic properties of projective and injective objects in this setting.

\begin{proposition}[{\cite[Lemma~3.3]{GKKP}}]\label{prop:right exact-is-enough-on-projectives}
  Let \(\cat{A}\) and \(\cat{B}\) be abelian categories,
  and assume that \(\cat{B}\) has enough projectives.
  Let \(\iota\from \proj{\cat{A}} \longhookrightarrow \cat{A}\) be the inclusion functor.
  Then the restriction functor \(\mathsf{Rex}(\cat{A}, \cat{B}) \xrightarrow{\blank \circ \iota} [\proj{\cat{A}}, \cat{B}]\) is an equivalence.
\end{proposition}

\begin{proposition}[{\cite[Lemma~3.3]{GKKP}}]\label{prop:left exact-is-enough-on-injectives}
  Let \(\cat{A}\) and \(\cat{B}\) be abelian categories,
  and assume that \(\cat{B}\) has enough injectives.
  Let \(\iota\from \inj{\cat{A}} \longhookrightarrow \cat{A}\) be the inclusion functor.
  Then \(\mathsf{Lex}(\cat{A},\cat{B}) \xrightarrow{\blank \circ \iota} [\inj{\cat{A}}, \cat{B}]\) is an equivalence.
\end{proposition}

We refer to~\cite[Chapter~1]{Etingof2015} for the following definitions and results.
An object of a \(\Bbbk\)\hyp{}linear category is said to be \emph{simple} if any non-zero endomorphism thereof is an isomorphism.
A category \(\cat{C}\) is said to be \emph{semisimple} if any of its objects is a direct sum of simple objects.

\begin{proposition}[{\cite[Theorem~C.6]{Turaev2017}}]\label{toomuchTV}
  A semisimple category is abelian.
  Any epimorphism or monomorphism in a semisimple category splits.
  In particular, any functor from a semisimple category to an abelian category is exact.
\end{proposition}

An object \(V\) in an abelian category \(\cat{A}\) has \emph{finite length} if there is a filtration
\[
  0 = V_0 \subseteq V_1 \subseteq \dots \subseteq V_n = V,
\]
such that \(V_i/V_{i-1}\) is simple, for all \(i \in \{\, 1, \dots, n \,\}\).

\begin{definition}\label{def:finite-abelian-category}
  A category \(\cat{C}\) is said to be \emph{finite abelian} if
  it is abelian,
  hom-finite\note{%
    A (\(\Bbbk\)\hyp{}linear) category \(\cat{C}\) is called \emph{hom-finite}
    if the hom space \(\cat{C}(x, y)\) is finite\hyp{}dimensional for all \(x, y \in \cat{C}\).%
  },
  has enough projectives,
  only finitely many isomorphism classes of simple objects, and
  all of its objects are of finite length.
\end{definition}

\begin{lemma}
  A category \(\cat{A}\) is finite abelian
  if and only if
  there is a finite\hyp{}dimensional \(\Bbbk\)\hyp{}algebra \(A\) such that \(\cat{A} \simeq \lmod{A}\).
\end{lemma}

\begin{definition}
  A category is said to be \emph{locally finite abelian} if
  it is abelian,
  hom-finite,
  and all of its objects are of finite length.
\end{definition}

\begin{lemma}\label{lem:campus-bar}
  A category \(\cat{C}\) is locally finite abelian
  if and only if
  there exists a \(\Bbbk\)\hyp{}coalgebra \(C\) such that \(\cat{C} \simeq \tetramodfd[C]{}\).
\end{lemma}

Locally finite abelian categories satisfy a Krull--Schmidt-type theorem.

\begin{proposition}[{\cite[Theorem~1.5.7]{Etingof2015}}]\label{prop:krull-schmidt}
  Every object of finite length admits a unique decomposition into a direct sum of indecomposable objects.
\end{proposition}

\begin{definition}\label{def:krull-schmidt-category}
  An abelian category \(\cat{C}\) that satisfies \cref{prop:krull-schmidt} is called a \emph{Krull--Schmidt category}.
\end{definition}

General abelian monoidal categories admit a very useful special case of \emph{Beck's monadicity theorem},
see for example~\cite[Section~4.1]{ben-zvi18:integ}.

\begin{theorem}\label{thm:SmallBeck}
  Let \(\adj{F}{U}{\cat{A}}{\cat{B}}\) be an adjunction of abelian categories.
  If \(U\) is right exact and reflects zero objects\note{%
    A functor \(F\) \emph{reflects zero objects} if \(Fx \cong 0 \implies x \cong 0\).%
  },
  the comparison functor \(K \from \cat{B} \to \cat{A}^{UF}\) is an equivalence.
  Likewise, if \(F\) is left exact and reflects zero objects
  then the comparison functor \(K\from \cat{A} \to \cat{B}^{FU}\) is an equivalence.
\end{theorem}

\subsection{Tensor and ring categories}

We refer the reader to~\cite[Chapter~4]{Etingof2015} for a comprehensive account of the theory of (multi)tensor and ring categories.
All categories and functors in this section are again assumed to be \(\Bbbk\)\hyp{}linear.

\begin{definition}\hfill
  \begin{itemize}
    \item A \emph{tensor category} is a locally finite abelian rigid monoidal category \(\cat{C}\) such that \(\mathrm{End}_{\cat{C}}(\1) \cong \Bbbk\).
    \item A \emph{finite tensor category} is a finite abelian rigid monoidal category \(\cat{C}\) such that \(\mathrm{End}_{\cat{C}}(\1) \cong \Bbbk\).
    \item A \emph{ring category} is a locally finite abelian separately exact monoidal category \(\cat{C}\) such that \(\mathrm{End}_{\cat{C}}(\1) \cong \Bbbk\).
    \item A \emph{finite ring category} is a finite abelian separately exact monoidal category \(\cat{C}\) such that \(\mathrm{End}_{\cat{C}}(\1) \cong \Bbbk\).
  \end{itemize}
\end{definition}

Despite apparent similarity in terminology,
the notion of a ring category is not related to that of a rig category,
such as those considered in~\cite{johnson-yau2021:TwoDimensionalCategories}.

\begin{definition}
  Let \(\cat{C}\) be a ring category.
  A \emph{fibre functor} for \(\cat{C}\) is a faithful and exact monoidal functor \(U\from \cat{C} \to \kvect\).
\end{definition}

The next result is a variant of Tannaka--Krein duality for ring categories.

\begin{theorem}[{\cite[Theorem~5.4.1]{Etingof2015}}]
  Let \(\cat{C}\) be a ring category, and assume it admits a fibre functor \(U\from \cat{C} \to  \kvect\).
  Then there exists a bialgebra \(B\) such that there is a monoidal equivalence \(K_B\from \cat{C} \to \lcomod{B}\)
  and a monoidal natural isomorphism \(U_{B} \circ K_{B} \cong U\), where \(U_{B}\from \lcomod{B} \to \kvect\) is the forgetful functor.
  Further, \(B\) is unique up to isomorphism, Hopf if and only if \(\cat{C}\) is a tensor category, and finite\hyp{}dimensional if and only if \(\cat{C}\) is a finite ring category.
\end{theorem}

\section{Algebra and module objects}\label{sec:algebra-and-module-objects}

\textsc{Let \(\cat{C}\) be a monoidal category}.
An \emph{algebra object in \(\cat{C}\)} consists of
an object \(A \in \cat{C}\)
together with a \emph{multiplication} \(\mu\from A \otimes A \to A\)
and a \emph{unit} \(\eta\from \1 \to A\),
satisfying associativity and unitality axioms analogous to those for a \(\Bbbk\)\hyp{}algebra.
In fact, a \(\Bbbk\)\hyp{}algebra is the same as an algebra object in \(\kVect\).

A \emph{coalgebra object} in \(\cat{C}\) is
an object \(C \in \cat{C}\) equipped
with a comultiplication \(\Delta\from C \to C \otimes C\)
and a counit morphism \(\varepsilon\from C \to \1\),
satisfying coassociativity and counitality axioms,
such that \(C\) becomes an algebra object in \(\cat{C}^{\op}\).

\begin{example}\label{ex:monoid-iff-tensoring-is-monad}
  An object \(A\) in a monoidal category \(\cat{C}\) is an algebra object in \(\cat{C}\)
  if and only if \(\blank \otimes A\from \cat{C}\to \cat{C}\) is a monad.
\end{example}

\begin{example}\label{ex:dual-pairs}
  Let \(\cat{C}\) be a monoidal category,
  and suppose that \(x \in \cat{C}\) has a left dual \(\ld{\!{}x} \in \cat{C}\).
  Then this induces an algebra object
  \[
    (x \otimes \ld{\!{}x},
    \ x \otimes \ev_x^{\ell} \otimes \ld{\!{}x}\from x \otimes \ld{\!{}x} \otimes x \otimes \ld{\!{}x} \to x \otimes \ld{\!{}x},
    \ \coev_x^{\ell} \from 1 \to x \otimes \ld{\!{}x})
  \]
  and, dually, a coalgebra object \((\ld{\!{}x} \otimes x,\, \ld{\!{}x} \otimes \coev_x^{\ell} \otimes X,\, \ev_x^{\ell})\) in \(\cat{C}\).
\end{example}

A special case of \cref{ex:dual-pairs} is when we consider the monoidal category \([\cat{C},\cat{C}]\) of endofunctors.
Then a left dual of some \(U \in [\cat{C},\cat{C}]\) is a left adjoint,
with \(\ev_U^{\ell}\) and \(\coev_U^{\ell}\) giving the coevaluation and evaluation, respectively.
the resulting algebra and coalgebra structures are the monad and comonad corresponding to the adjunction.

\begin{definition}\label{ModuleInModule}
  Let \((\cat{M}, \lact)\) be a left \(\cat{C}\)\hyp{}module category
  and \((A, \mu, \eta)\) an algebra object in \(\cat{C}\).
  A \emph{left \(A\)\hyp{}module in \(\cat{M}\)} is
  an object \(M \in \cat{M}\)
  together with an \emph{action} morphism \(\alpha\from A \lact M \to M\)
  such that the following diagrams commute
  \[
    \begin{mytikzcd}[ampersand replacement=\&]
      {(A \otimes A) \triangleright M} \& {A \triangleright (A \triangleright M)} \& {A \triangleright M} \& {\1 \lact M} \& M \\
      {A \triangleright M} \&\& M \& {A \lact M}
      \arrow["{{{\cat{M}_{\mathsf{a}}}}}", from=1-1, to=1-2]
      \arrow["{{\mu \,\triangleright\, m}}"', from=1-1, to=2-1]
      \arrow["{{A \,\triangleright\, \alpha}}", from=1-2, to=1-3]
      \arrow["{{\alpha}}", from=1-3, to=2-3]
      \arrow["{{\cat{M}_{\mathsf{u}}}}", from=1-4, to=1-5]
      \arrow["{\eta \,\lact\, M}"', from=1-4, to=2-4]
      \arrow["{{\alpha}}"', from=2-1, to=2-3]
      \arrow["{\alpha}"', from=2-4, to=1-5]
    \end{mytikzcd}
  \]

  A \emph{morphism of modules} is an arrow \(f\from M \to N\) in \(\cat{M}\) commuting with the respective actions
  in the sense that \(f \circ \alpha_m = \alpha_n \circ (A \lact f)\).
\end{definition}

\begin{example}\label{MoreModules}
  Clearly, left \(A\)\hyp{}modules in \(\cat{M}\) and their morphisms form a category,
  which we shall denote by \(\modma\).
  The obvious forgetful functor
  \(\forgetleftmod{A}{\cat{M}}\from \modma \to \cat{M}\)
  admits a left adjoint
  \begin{align*}
    \freeleftmod{A}{\cat{M}}\from \cat{M} &\to \leftmod{A}{\cat{M}} \\[-0.25cm]
    M &\mapsto \big(A \triangleright M,\ A \lact (A \lact M) \xrightarrow{\cat{M}_{\mathrm{a}}^{-1}} (A \otimes A) \lact M \xrightarrow{\mu \,\triangleright\, M} A \lact M\big)
  \end{align*}
  The unit of this adjunction is given by \(M \xrightarrow{\,\cat{M}_{\mathrm{u}}\,} \1 \triangleright M \xrightarrow{\eta \,\triangleright\, M} A \triangleright M\),
  and the counit is \(\alpha_M \from A \triangleright M \to M\).
\end{example}

Given a coalgebra object \(C\) in \(\cat{C}\),
a \emph{left \(C\)\hyp{}comodule in \(\cat{M}\)} is an object \(N \in \cat{C}\) satisfying axioms dual to those for a module in \(\cat{M}\),
so that a left \(C\)\hyp{}comodule in \(\cat{M}\) is the same as a left \(C\)\hyp{}module in the \(\cat{C}^{\op}\)\hyp{}module category \(\cat{M}^{\op}\).
Comodule morphisms are defined similarly.
We denote the category of \(C\)\hyp{}comodules in \(\cat{M}\) by \(\comodmc\).
There is an adjunction analogous to that for modules:
\[
  \adj{\forgetleftcomod{C}{\cat{M}}}{\freeleftcomod{C}{\cat{M}}}{\cat{M}}{\comodmc}.
\]

Similarly, one can define right modules and comodules in a right module category,
and bimodules and bicomodules in a bimodule category.

\begin{definition}\label{def:comodule-algebra}
  Let \(A\) be an algebra in \(\cat{C}\).
  A (left) \emph{\(A\)\hyp{}comodule algebra} is an algebra \(C\), together with a left \(A\)\hyp{}comodule structure \(\rho\from C \to A \otimes C\),
  such that \(\rho\) is a morphism of \(A\)\hyp{}algebras.
\end{definition}

Analogously to \cref{def:comodule-algebra} one defines right comodule algebras, as well as comodule coalgebras and module algebras.

\begin{notation}\label{notation-for-vect-modules}
  If the underlying category is the category of vector spaces,
  then we use the following notation,
  to differentiate it from the general case:
  \[
    \Tetramod[C]{} \defeq \leftcomod{C}{\kVect}.
  \]
  for the category of left \(C\)\hyp{}comodules of a coalgebra \(C\).

  Analogously, we for example define \(\Tetramod[][A]{}\), \(\Tetramod[C][][A][]\);
  or \(\tetramodfd[C]{}\), \(\tetramodfd[][A]{}\),
  and \(\tetramodfd[B][B][B][B]\) for the finite\hyp{}dimensional case.
\end{notation}

\begin{remark}\label{rem:lifting-functors-to-mod}
  Let \(A\) be an algebra object in \(\mathcal{C}\) and let \(\mathcal{M}\) and \(\cat{N}\) be left \(\mathcal{C}\)\hyp{}module categories.
  A lax \(\mathcal{C}\)\hyp{}module functor \(F\from \mathcal{N} \to \mathcal{M}\)
  can be lifted to
  \[
    \leftmod{A}{F}\from \leftmod{A}{\mathcal{N}} \to \leftmod{A}{\mathcal{M}}
  \]
  on the respective categories of \(A\)\hyp{}modules
  by sending \(N \in \leftmod{A}{\mathcal{N}}\) to the left \(A\)\hyp{}module \(FN\),
  whose action is given by
  \[
    \alpha_{FN}\from A \lact_{\mathcal{M}} FN \xrightarrow{F_{\mathrm{a};A,N}} F(A \lact_{\mathcal{N}} N) \xrightarrow{F \alpha_{N}} FN.
  \]
  This functor is a lift of \(F\)—it satisfies the following relation:
  \[
    \begin{mytikzcd}[ampersand replacement=\&]
      {\leftmod{A}{\mathcal{N}}} \&\& {\leftmod{A}{\mathcal{M}}} \\
      {\mathcal{N}} \&\& {\mathcal{M}}
      \arrow["{\leftmod{A}{F}}", from=1-1, to=1-3]
      \arrow["{\forgetleftmod{A}{\mathcal{N}}}"', from=1-1, to=2-1]
      \arrow["{\forgetleftmod{A}{\mathcal{M}}}", from=1-3, to=2-3]
      \arrow["F"', from=2-1, to=2-3]
    \end{mytikzcd}
  \]
\end{remark}

\begin{example}\label{ex:bicategory-of-bimodules}
  Let \(\cat{C}\) be a cocomplete symmetric monoidal category,
  such that its is right exact in both variables.
  Then one may define the bicategory \(\BiCat{Bimod}(\cat{C})\) of bimodules in \(\cat{C}\) as follows:
  \begin{itemize}
    \item Objects are algebra objects in \(\cat{C}\).
    \item For \(A\) and \(B\) in \(\cat{C}\), the hom-category \(\BiCat{Bimod}(\cat{C})(A, B)\)
    is given by \(B\)\hyp{}\(A\)\hyp{}bimodules in \(\cat{C}\) and their homomorphisms.
    Horizontal composition is given by the balanced tensor product of bimodules.
  \end{itemize}
\end{example}

Analogously to \cref{ex:bicategory-of-bimodules},
there is a bicategory \(\BiCat{Bicomod}(\cat{C})\) of \(\cat{C}\)\hyp{}bicomodules,
for \(\cat{C}\) a complete symmetric monoidal category,
with the tensor product being left exact in both variables.

\begin{definition}\label{def:filtered-category}
  A category \(\cat{I}\) is \emph{filtered} if every finite diagram has a cocone.
\end{definition}

Equivalently, \cref{def:filtered-category} could be stated in the following way:
a category \(\cat{I}\) is filtered if the following two conditions hold:
\begin{itemize}
  \item for all \(x,y \in \cat{I}\)
  there exist arrows \(k_x\from x \to k\) and \(k_y\from y \to k\) in \(\cat{I}\); and
  \item for all parallel arrows \(f,g \from x \to y\) in \(\cat{I}\)
  there exists a \(k_{fg}\from y \to k\), such that \(k_{fg} \circ g = k_{fg} \circ f\).
\end{itemize}

\begin{definition}\label{def:filtered-colimit}
  A diagram \(F\from \cat{I} \to \cat{C}\) is called \emph{filtered} if the domain \(\cat{I}\) is a filtered category.
  A \emph{filtered colimit} is a colimit of a filtered diagram,
  and a functor is called \emph{finitary} if it preserves filtered colimits.
\end{definition}

Ordinarily,
limits and colimits do not commute.
For example, for sets \(X\), \(Y\), \(W\), and \(V\),
the canonical morphism
\((X \times Y) \coprod\, (W \times V) \to (X \coprod V) \times (Y \coprod V)\)
is not isomorphism in general.
However, it turns out that filtered colimits do commute with finite limits,
see for example~\cite[Theorem~\textsc{ix}.2.1]{MacLane1998}.

\begin{proposition}\label{prop:filtered-colimits-and-finite-limits-commute}
  Filtered colimits commute with finite limits in \(\Set\):
  given a functor \(F\from \cat{J} \times \cat{I} \to \Set\) of two diagram categories,
  where \(\cat{J}\) is finite and \(\cat{I}\) is small and filtered,
  the following canonical morphism is an isomorphism:
  \[
    \colim\displaylimits_i \lim_j F(j, i) \iso \lim_j \colim\displaylimits_i F(j, i).
  \]
\end{proposition}

\begin{lemma}\label{finitarythings}
  Let \(\cat{A}\) and \(\cat{B}\) be additive categories.
  Finite biproducts endow \(\cat{A}\) and \(\cat{B}\) with the structure of \(\kvect\)\hyp{}module categories.
  Any functor \(F\from \cat{A} \to \cat{B}\) is a \(\kvect\)\hyp{}module functor with respect to these \(\kvect\)\hyp{}actions.

  If \(\cat{B}\) admits filtered colimits,
  \(F\) extends essentially uniquely to a finitary \(\kVect\)\hyp{}module functor \(\Ind(\cat{A}) \to \cat{B}\),
  and any finitary \(\kVect\)\hyp{}module functor is of this form.
\end{lemma}
\begin{proof}
  Let \(\mathsf{add}(\cat{A})\) denote the cocompletion of \(\cat{A}\) under finite biproducts—the \emph{additive closure} of \(\cat{A}\).
  As finite biproducts are absolute colimits,
  the left adjoint \(\mathsf{add}(\cat{A}) \to \cat{A}\) of the canonical inclusion \(\cat{A} \longhookrightarrow \mathsf{add}(\cat{A})\),
  coming from the universal property of \cref{cocompletingmodules},
  is an equivalence of \(\kvect\)\hyp{}module categories.
  Hence, so is the inclusion \(\cat{A} \longhookrightarrow \mathsf{add}(\cat{A})\).
  We find the equivalences
  \[
    \Cat_{\Bbbk}(\cat{A},\blank)
    \simeq \lMod{\kvect}(\mathsf{add}(\cat{A}),\blank)
    \simeq \lMod{\kvect}(\cat{A},\blank),
  \]
  establishing the first claim.
  The second claim follows in a similar way,
  using that \(\kVect \simeq \Ind(\kvect)\) and
  \[
    \lMod{\kvect}(\cat{A},\cat{B})
    \simeq \lMod{\kVect}_{\mathrm{filt}}(\Ind(\cat{A}), \cat{B}).\qedhere
  \]
\end{proof}

\begin{example}\label{ex:cotensor-product-of-comodules}
  Let \(C\) be a coalgebra in \(\kVect\).
  Given two \(C\)\hyp{}comodules \((X, \rho)\) and \((Y, \lambda)\),
  one may form their \emph{cotensor product}—the vector space \(X \cotens Y\) given by the equaliser
  \[
    \begin{mytikzcd}[ampersand replacement=\&]
      {X \cotens Y} \& {X \kotimes Y} \& {X \kotimes C \kotimes Y.}
      \arrow[hook, from=1-1, to=1-2]
      \arrow["{\rho \kotimes Y}", shift left=1, from=1-2, to=1-3]
      \arrow["{X \kotimes \lambda}"', shift right=1, from=1-2, to=1-3]
    \end{mytikzcd}
  \]
\end{example}

Sometimes, the formalism of module categories%
—even those over the seemingly trivial monoidal categories \(\kvect\) and \(\kVect\)—%
provides additional clarity to classical statements about algebraic and coalgebraic \(\Bbbk\)\hyp{}linear structures.
We give a brief proof of the result below as an example of this.

\begin{proposition}[{\cite[Proposition~2.1]{takeuchi77:morit}}]\label{TakeuchiEilenbergWatts}
  Let \(\cat{D}\) be an abelian category admitting filtered colimits, and let \(C\) be a coalgebra over \(\Bbbk\).
  There is an equivalence
  \[
    \begin{aligned}
      \mathsf{Lexf}(\Tetramod[C]{}, \cat{D}) &\iso \rightcomod{\cat{D}}{C} \\
      F &\mapsto FC \\
      N \cotens \blank &\longmapsfrom N,
    \end{aligned}
  \]
  where the left-hand side is the category of left exact finitary functors from \(\Tetramod[C]{}\),
  and the right-hand side—following the notation of \cref{MoreModules}—denotes the category of right \(C\)\hyp{}comodules in \(\cat{D}\).
  Here, \(\cat{D}\) is endowed with the \(\kVect\)\hyp{}module structure described in \cref{FinitaryAction}.

  In particular, for a coalgebra \(D\) over \(\Bbbk\), we have
  \[
    \mathsf{Lexf}(\Tetramod[C]{}, \Tetramod[D]{}) \simeq \Tetramod[D][][][C].
  \]
\end{proposition}
\begin{proof}
  Since \(\Tetramod[C]{} \simeq \Ind(\tetramodfd[C]{})\),
  the functor \(F\) can be seen as induced to filtered colimits from its restriction to \(\tetramodfd[C]{}\),
  and as such is a \(\kVect\)\hyp{}module functor, following \cref{finitarythings}.
  Observe that \(C\), being a bicomodule over itself, is a right \(C\)\hyp{}comodule in \(\Tetramod[C]{}\).
  Thus, \(FC\) is a right \(C\)\hyp{}comodule in \(\cat{D}\).

  To see that \(F \cong FC \cotens \blank\), let \(N \in \Tetramod[C]{}\) and recall that the bar construction provides a functorial injective resolution
  \[
    N \cong \eq(\!\!
    \begin{mytikzcd}[ampersand replacement=\&,sep=small]
      {C \otimes N} \& { C \otimes C \otimes N}
      \arrow["{\Delta_{C} \otimes N}", shift left, from=1-1, to=1-2]
      \arrow["{C \otimes \Delta_{N}}"', shift right, from=1-1, to=1-2]
    \end{mytikzcd}
    \!\!).
  \]
  Thus,
  \[
    \begin{aligned}
      FN
      &\cong F\big(\eq(\!\!
        \begin{mytikzcd}[ampersand replacement=\&,sep=small]
          {C \otimes N} \& { C \otimes C \otimes N}
          \arrow["{\Delta_{C} \otimes N}", shift left, from=1-1, to=1-2]
          \arrow["{C \otimes \Delta_{N}}"', shift right, from=1-1, to=1-2]
        \end{mytikzcd}
        \!\!)\big)
        \cong \eq\big(\!\!
        \begin{mytikzcd}[ampersand replacement=\&,sep=small]
          {F(C \otimes N)} \& {F(C \otimes C \otimes N)}
          \arrow["{F(\Delta_{C} \otimes N)}", shift left, from=1-1, to=1-2]
          \arrow["{F(C \otimes \Delta_{N})}"', shift right, from=1-1, to=1-2]
        \end{mytikzcd}
        \!\!\big) \\
      &\cong \eq\big(\!\!
        \begin{mytikzcd}[ampersand replacement=\&,sep=small]
          {FC \otimes N} \& {FC \otimes C \otimes N}
          \arrow["{F(\Delta_{C}) \otimes N}", shift left, from=1-1, to=1-2]
          \arrow["{FC \otimes \Delta_{N}}"', shift right, from=1-1, to=1-2]
        \end{mytikzcd}
        \!\!\big)
        = FC \cotens N,
    \end{aligned}
  \]
  where the first isomorphism is using the bar construction,
  the second follows from left exactness of \(F\),
  the third uses the fact that \(F\) is a right \(\kVect\)\hyp{}module functor,
  and the fourth follows by observing that \(F(\Delta_{C}) = \Delta_{F(C)}\).
\end{proof}

\section{Monoidal bicategories}\label{sec:monoidal-bicategories}

\textsc{We briefly recall the notion of a monoidal bicategory},
which will play an important role in the string diagrams used throughout.
More detailed accounts are given~\cite[Chatper~12]{johnson-yau2021:TwoDimensionalCategories},
see also~\cite{gordon95:coher,gurski2006algebraic,garner16:enric}.

\begin{definition}
  A \emph{monoidal bicategory} consists of
  a bicategory \(\mathbb{B}\) together with the following data:
  \begin{itemize}
    \item A pseudofunctor \(\boxtimes\colon \mathbb{B} \times \mathbb{B} \to \mathbb{B}\).
    \item A pseudofunctor \(\1_{\mathbb{B}}\colon \terminal \to \mathbb{B}\).
    The image of the unique object \(\terminal\) in \(\mathbb{B}\)
    will be referred to as the \emph{identity object} of \(\mathbb{B}\),
    and, abusing notation,
    also denoted by \(\1_{\mathbb{B}}\).
    The pseudofunctoriality of \(\1_{\mathbb{B}}\) yields additional structure:
    a 1-morphism
    \(\mathrm{I}_{\mathbb{B}}\colon \1_{\mathbb{B}} \to \1_{\mathbb{B}}\)
    and invertible 2-morphisms
    \(\nabla^{\mathrm{I}}\colon \mathrm{I}_{\mathbb{B}} \circ \mathrm{I}_{\mathbb{B}} \iso \mathrm{I}_{\mathbb{B}}\)
    and \(\eta^{\mathrm{I}}\colon \Id_{\1_{\mathbb{B}}} \iso \mathrm{I}_{\mathbb{B}}\),
    such that \((\1_{\mathbb{B}},\nabla^{\mathrm{I}},\eta^{\mathrm{I}})\)
    form an idempotent monad in \(\mathbb{B}\).
    \item An adjoint pseudonatural equivalence
    \((\mathfrak{A}, \mathfrak{A}^{\blackdiamond}, \eta^{\mathfrak{A}}, \varepsilon^{\mathfrak{A}})\):
    \[
      \begin{mytikzcd}[ampersand replacement=\&]
	{\mathbb{B} \times \mathbb{B} \times \mathbb{B}} \& {\mathbb{B} \times \mathbb{B}}\\
	{\mathbb{B} \times \mathbb{B}} \& {\mathbb{B}}
	\arrow["{\boxtimes \times \mathbb{B}}", from=1-1, to=1-2]
	\arrow["{\mathbb{B} \times \boxtimes}"', from=1-1, to=2-1]
	\arrow["{\mathfrak{A}}", shorten <=12pt, shorten >=12pt, Rightarrow, from=1-2, to=2-1]
	\arrow["\boxtimes", from=1-2, to=2-2]
	\arrow["\boxtimes"', from=2-1, to=2-2]
      \end{mytikzcd}
    \]
    which will be referred to as
    the \emph{associator} for \(\mathbb{B}\).
    Its components are 1-morphisms of the form
    \(\mathfrak{A}_{\cat{A,B,C}}\colon (\cat{A} \boxtimes \cat{B}) \boxtimes \cat{C} \iso \cat{A} \boxtimes (\cat{B} \boxtimes \cat{C})\).
    \item Pseudonatural adjoint equivalences
    \((\mathfrak{L},\mathfrak{L}^{\blackdiamond}, \eta^{\mathfrak{L}}, \varepsilon^{\mathfrak{L}})\)
    and \((\mathfrak{R},\mathfrak{R}^{\blackdiamond}, \eta^{\mathfrak{R}}, \varepsilon^{\mathfrak{R}})\),
    referred to as \emph{left and right unitors}, respectively:
    \[
      \begin{mytikzcd}[ampersand replacement=\&,sep=scriptsize]
        {\mathbb{B}} \&\& {\mathbb{B} \times \mathbb{B}} \& {\mathbb{B}}
        \arrow["\boxtimes", from=1-3, to=1-4]
        \arrow["{\id_{\mathbb{B}}\times \1_{\mathbb{B}}}"{pos=0.7}, shift left=2, from=1-1, to=1-3]
        \arrow["{ \1_{\mathbb{B}}\times \id_{\mathbb{B}}}"'{pos=0.7}, shift right=2, from=1-1, to=1-3]
        \arrow[""{name=0, anchor=center, inner sep=0}, shift left=2, curve={height=-30pt}, from=1-1, to=1-4]
        \arrow[""{name=1, anchor=center, inner sep=0}, shift right=2, curve={height=30pt}, from=1-1, to=1-4]
        \arrow["{\mathfrak{L}}"', shorten <=4pt, Rightarrow, from=0, to=1-3]
        \arrow["{\mathfrak{R}}"', shorten <=4pt, Rightarrow, from=1, to=1-3]
      \end{mytikzcd}
    \]
    Its components are 1-morphisms of \(\mathbb{B}\)
    of the form \(\mathfrak{R}\colon \1_{\mathbb{B}} \boxtimes \cat{A} \iso \cat{A}\)
    and \(\mathfrak{L}\colon \cat{A} \boxtimes \1_{\mathbb{B}} \iso \cat{A}\),
    for \(\cat{A} \in \Ob\mathbb{B}\).
    \item An invertible modification \(\mathfrak{p}\),
    as indicated by the components of \(\mathfrak{p}\) displayed below,
    where \(\cat{A},\cat{B},\cat{C},\cat{D} \in \Ob\mathbb{B}\):
    \begin{equation}\label[diagram]{Pentagonator}
      \begin{mytikzcd}[ampersand replacement=\&, column sep=tiny]
        {(\cat{D} \boxtimes (\cat{C} \boxtimes \cat{B}) ) \boxtimes \cat{A}} \&\& {\cat{D} \boxtimes ((\cat{C} \boxtimes \cat{B}) \boxtimes \cat{A})} \\
        {((\cat{D} \boxtimes \cat{C}) \boxtimes \cat{B}) \boxtimes \cat{A}} \&\& {\cat{D} \boxtimes (\cat{C} \boxtimes (\cat{B} \boxtimes \cat{A}))} \\
        \& {(\cat{D} \boxtimes \cat{C}) \boxtimes (\cat{B} \boxtimes \cat{A})}
        \arrow[""{name=0, anchor=center, inner sep=0}, from=1-1, to=1-3]
        \arrow[from=2-1, to=1-1]
        \arrow[from=1-3, to=2-3]
        \arrow[from=2-1, to=3-2]
        \arrow[from=2-3, to=3-2]
        \arrow["{\mathfrak{p}_{\cat{D},\cat{C},\cat{B},\cat{A}}}", shorten <=16pt, shorten >=16pt, Rightarrow, from=0, to=3-2]
      \end{mytikzcd}
    \end{equation}
    The arrows in the diagram are formed using the associator \(\mathfrak{A}\) as indicated. We refer to \(\mathfrak{p}\) as the \emph{pentagonator}.
    \item Invertible modifications \(\mathfrak{l}\), \(\mathfrak{m}\), and \(\mathfrak{r}\);
    the \emph{left, middle, and right 2-unitors}:
    \[
      \mathscale{0.9}{\begin{mytikzcd}[ampersand replacement=\&]
	{(\cat{B} \boxtimes \1_{\mathbb{B}}) \boxtimes \cat{A}} \& {\cat{B} \boxtimes (\1_{\mathbb{B}} \boxtimes \cat{A})} \& {(\cat{B} \boxtimes \cat{A}) \boxtimes \1_{\mathbb{B}}} \& {\cat{B} \boxtimes (\cat{A} \boxtimes \1_{\mathbb{B}})} \\
	{\cat{B} \boxtimes \cat{A}} \& {\cat{B} \boxtimes \cat{A}} \& {\cat{B} \boxtimes \cat{A}} \\
	\& {(\1_{\mathbb{B}} \boxtimes \cat{B}) \boxtimes \cat{A}} \& {\1_{\mathbb{B}} \boxtimes (\cat{B} \boxtimes \cat{A})} \\
	\& {\cat{B} \boxtimes \cat{A}}
	\arrow[""{name=0, anchor=center, inner sep=0}, "{{\mathfrak{A}_{\cat{B},\1_{\mathbb{B}},\cat{A}}}}", from=1-1, to=1-2]
	\arrow["{{\cat{B} \boxtimes \mathfrak{L}_{\cat{A}}}}", from=1-2, to=2-2]
	\arrow["{{\mathfrak{A}_{\cat{B},\cat{A},\1_{\mathbb{B}}}}}", from=1-3, to=1-4]
	\arrow["{{\mathfrak{R}^{\blackdiamond}\boxtimes \cat{A}}}", from=2-1, to=1-1]
	\arrow[""{name=1, anchor=center, inner sep=0}, equals, from=2-1, to=2-2]
	\arrow["{{\mathfrak{R}_{\cat{B}\boxtimes \cat{A}}^{\blackdiamond}}}", from=2-3, to=1-3]
	\arrow[""{name=2, anchor=center, inner sep=0}, "{{\cat{B}\boxtimes \mathfrak{R}_{\cat{A}}^{\blackdiamond}}}"', curve={height=12pt}, from=2-3, to=1-4]
	\arrow["{{\mathfrak{A}_{\1_{\mathbb{B}},\cat{B},\cat{A}}}}", from=3-2, to=3-3]
	\arrow["{{\mathfrak{L}_{\cat{B}}\boxtimes \cat{A}}}"', from=3-2, to=4-2]
	\arrow[""{name=3, anchor=center, inner sep=0}, "{{\mathfrak{L}_{\cat{B}\boxtimes \cat{A}}}}", curve={height=-12pt}, from=3-3, to=4-2]
	\arrow["{\mathfrak{m}_{\cat{B},\cat{A}}}", shorten <=12pt, shorten >=12pt, Rightarrow, from=0, to=1]
	\arrow["{\mathfrak{r}_{\cat{B},\cat{A}}}"{pos=0.8}, shorten <=16pt, shorten >=8pt, Rightarrow, from=2, to=1-3]
	\arrow["{\mathfrak{l}_{\cat{B},\cat{A}}}"'{pos=0.2}, shorten <=8pt, shorten >=16pt, Rightarrow, from=3-2, to=3]
      \end{mytikzcd}}
    \]
  \end{itemize}

  \noindent{}This data is subject to coherence axioms,
  see~\cite[Explanation~12.1.3]{johnson-yau2021:TwoDimensionalCategories}.
\end{definition}

\begin{example}\label{MonoidalBicatCat}
  The 2-category \(\BiCat{Cat}_{\Bbbk}\) of
  \(\Bbbk\)\hyp{}linear categories,
  \(\Bbbk\)\hyp{}linear functors,
  and natural transformations between them
  is a monoidal bicategory
  when endowed with the \emph{tensor product} of \(\Bbbk\)\hyp{}linear categories:
  given \(\cat{A}\) and \(\cat{B}\), define the \(\Bbbk\)\hyp{}linear category \(\cat{A} \otimes_{\Bbbk} \cat{B}\) by
  \begin{itemize}
    \item \(\Ob(\cat{A} \otimes_{\Bbbk} \cat{B}) \defeq \Ob\cat{A} \times \Ob\cat{B}\),
    \item \((\cat{A}\otimes_{\Bbbk} \cat{B})((a,b),(a',b')) \defeq \cat{A}(a,a') \otimes_{\Bbbk} \cat{B}(b,b')\).
  \end{itemize}
  Given \(\Bbbk\)\hyp{}linear functors \(F\colon \cat{A} \rightarrow \cat{B}\) and \(G\colon \cat{C} \rightarrow \cat{D}\),
  the tensor product \(F \otimes_{\Bbbk} G\) is defined by
  \(\Ob(F\otimes_{\Bbbk} G) = \Ob(F) \times \Ob(G)\) and
  \begin{align*}
    {(F\otimes_{\Bbbk} G)}_{(a, b),(c, d)} \from \cat{A}(a, c) \otimes_{\Bbbk} \cat{B}(b, d) \xrightarrow{F_{a,c} \otimes_{\Bbbk} G_{b,d}} \cat{C}(Fa, Fc) \otimes_{\Bbbk} \cat{D}(Gb, Gd).
  \end{align*}
  The tensor product of natural transformations extends similarly.

  The associators and unitors are lifted from \(\Set\)
  and \(\kVect\),
  and the pentagonators and \(2\)\hyp{}unitors are trivial.
  For example,
  given \(\Bbbk\)\hyp{}linear categories \(\cat{A}\), \(\cat{B}\), and \(\cat{C}\),
  the associator
  \(\mathfrak{A}_{\cat{A}, \cat{B},\cat{C}}\)
  is the equivalence
  \begin{align*}
    (\cat{A} \otimes_{\Bbbk} \cat{B}) \otimes_{\Bbbk} \cat{C} &\iso \cat{A} \otimes_{\Bbbk} (\cat{B} \otimes_{\Bbbk} \cat{C}) \\
    ((a,b),c) &\mapsto (a,(b,c)) \\
    (f \otimes_{\Bbbk} g) \otimes_{\Bbbk} h &\mapsto f \otimes_{\Bbbk} (g \otimes_{\Bbbk} h).
  \end{align*}
  \Cref{Pentagonator} commutes in this case,
  and so the pentagonator may be chosen to be the identity.
\end{example}

Generalising the above example,
for a symmetric closed monoidal category \((\cat{V}, \otimes)\),
the bicategory \(\VCat\) is monoidal,
with the monoidal structure \(\cat{A} \otimes \cat{B}\) defined analogously to \(\cat{A} \otimes_{\Bbbk} \cat{B}\) in the case of \(\cat{V} \defeq \kVect\).

\begin{example}\label{ex:bimod}
  The bicategory \(\BiCat{Bimod}(\cat{C})\) of \cref{ex:bicategory-of-bimodules}
  is a monoidal bicategory with respect to the underlying tensor product of \(\cat{C}\).
  Tensoring bimodules \({}_{B}M_{A}\) and \({}_{D}N_{C}\) is given by
  the \((B \otimes D)\)\hyp{}\((A \otimes C)\)\hyp{}bimodule
  \({}_{B}M_{A} \otimes {}_{D}N_{C}\),
  and this extends to bimodule homomorphisms in the obvious way.
\end{example}

% The coherence structure of the example monoidal bicategories that follow are not of great interest to us, so, in contrast to \cref{MonoidalBicatCat}, where we specified it for completeness, we will omit it in the remaining cases. We will however provide references to articles where more details can be found.

For details on the following example,
together with its more general variant where \(\cat{V}\) is only required to be braided,
see~\cite[Section~7]{day97:monoid-hopf}.

\begin{example}\label{MonoidalBicatProf}
  The bicategory \(\Prof_{\Bbbk}\) has as objects \(\Bbbk\)\hyp{}linear categories,
  and the hom-categories are given by
  \[
    \Prof_{\Bbbk}(\cat{C}, \cat{D}) \defeq \BiCat{Cat}_{\Bbbk}(\cat{D}^{\op} \kotimes \cat{C}, \kVect).
  \]
  We call \(F \in \Prof_{\Bbbk}(\cat{C}, \cat{D})\) a \emph{profunctor},
  and write \(F \from \cat{C} \slashedrightarrow \cat{D}\).
  Composition of profunctors \(F\from \cat{C} \slashedrightarrow \cat{D}\) and \(G\from \cat{D} \slashedrightarrow \cat{E}\) is given by the coend
  \[
    (G \diamond F)(e, c) \defeq \int^{d \in \cat{D}} G(e, d) \kotimes F(d, c).
  \]

  Further, \(\Prof_{\Bbbk}\) is monoidal when endowed the monoidal structure \(\otimes_{\Bbbk}\).
  On objects, it coincides with that in \cref{MonoidalBicatCat}:
  given categories \(\cat{A}\) and \(\cat{B}\),
  their tensor product is \(\cat{A} \otimes_{\Bbbk} \cat{B}\).
  Given profunctors \(\Phi\colon \cat{A} \slashedrightarrow \cat{A}'\) and \(\Psi\colon \cat{B} \slashedrightarrow \cat{B}'\),
  the profunctor \(\Phi \otimes_{\Bbbk} \Psi\colon \cat{A} \otimes_{\Bbbk} \cat{A}' \slashedrightarrow \cat{B} \otimes_{\Bbbk} \cat{B}'\)
  is defined via
  \[
    (\Phi \otimes_{\Bbbk} \Psi)((b, b'),(a, a')) = \Phi(a', a) \otimes_{\Bbbk} \Psi(b', b).
  \]
  This extends similarly to natural transformations between profunctors.
\end{example}

\subsection{String diagrams in monoidal bicategories}\label{sec:string-diagr-mono}

\textsc{This section extends the graphical language} of
\cref{sec:graphical-calculus-introduction}
by incorporating the tensor product of \(\BiCat{Cat}\),
which turns it into a monoidal 2-category.
In our presentation, we closely follow Willerton~\cite{Willerton2008},
see~\cite{day03:lax,street03:funct,barrett24:gray,demirdilek25:surfac-diagr-froben-algeb-froben}
for other examples of ``sheet'' or ``surface'' diagrams.

As before, we consider strings and vertices between them.
These are labelled with functors and natural transformations, respectively.
The strings and vertices are embedded into bounded rectangles, which we will call sheets.
Each (connected) region of a sheet is decorated with a category.
The same mechanics as for ordinary 2-dimensional string diagrams apply:
horizontal and vertical gluing represents composition of functors and natural transformations.
On top of these operations,
we add stacking sheets behind each other to depict the monoidal product of \(\BiCat{Cat}\).
Our convention is to read diagrams from front to back,
right to left, and bottom to top.

Given a monoidal category \(\cat{C}\),
two of the most vital building blocks in this new graphical language are
its tensor product \(\otimes\from \cat{C} \times \cat{C} \to \cat{C}\) and unit \(1 \in \cat{C}\);
they are shown in~\cref{fig:tensor-prod-unit-monoidal-cat}.
\begin{figure}[htbp]
  \centering
  \tikzfig{tensor-product}
  \scaption[-1.7cm]{%
    Graphical representation of the tensor product and unit of a monoidal category \((\cat{C}, \otimes, 1)\).%
    \label{fig:tensor-prod-unit-monoidal-cat}%
  }
\end{figure}
On the left, there are two sheets equating to two copies of \(\cat{C}\)
joined by a line: the tensor product of \(\cat{C}\).
On the right,
we have the unit of \(\cat{C}\)
considered as a functor
from the terminal category to \(\cat{C}\).
We represent \(\1\) by the empty sheet,
and the unit of \(\cat{C}\) by a dashed line.

\begin{example}\label{ex:graphical-opmonoidal-functor}
  Consider an oplax monoidal functor \((F, F_2, F_0) \from \cat{C} \to \cat{D}\).
  \Cref{fig:opmonoidal-comult-graphically} depicts the comultiplication \(F_2\) as a ``time evolution'',
  where we start with \(F(\blank \otimes \bblank)\) on the bottom,
  and end up with \(F(\blank) \otimes F(\bblank)\) at the top.
  The coassociativity and counitality conditions are depicted in \cref{fig:coassoc-and-counital-opmonoidal-functor}.
  \begin{figure}[htbp]
    \centering
    \tikzfig{opmonoidal-functor-delta-graphically}
    \vspace{\onelineskip}\hrule
    \scaption[-4.5cm]{%
      The comultiplication
      of an oplax monoidal functor \(F\from \cat{C} \to \cat{D}\).%
      \label{fig:opmonoidal-comult-graphically}%
    }
  \end{figure}
  \begin{figure}[htbp]
    \centering
    \tikzfig{coassoc-and-counital-opmonoidal-functor}\vspace{-\onelineskip}
    \scaption[-1cm]{%
      Coassociativity and counitality conditions of an oplax monoidal functor.%
      \label{fig:coassoc-and-counital-opmonoidal-functor}%
    }
  \end{figure}
\end{example}

\newpage% XXX: Manual orphan prevention. This results in a pretty empty
        % page, but since we have two large figures at the top of the
        % next one, it seems like the lesser of the two evils.
\begin{remark}
  Given two oplax monoidal functors \(F,G\from \cat{C} \to \cat{D}\),
  the string diagrammatic conditions for a natural transformation \(\eta\from F \nt G\) to be one of oplax monoidal functors
  look like
  \[
    \tikzfig{nat-of-oplax-monoidal-fun}
  \]
\end{remark}

\begin{example}\label{ex:graphical-opmonoidal-adjunctions}
  Suppose that \(\stdadj\) is an oplax monoidal adjunction.
  In string diagrams, the conditions that the unit \(\eta\) and the counit \(\varepsilon\)
  are oplax monoidal natural transformation is
  given in \cref{fig:comonoidal-adjunction-epsilon}.
  \begin{figure}[htbp]
    \centering
    \tikzfig{comonoidal-adjunction-epsilon}\vspace{-0.2cm}% XXX
    \scaption[-1cm]{%
      Conditions for \(\eta\) and \(\varepsilon\) to be oplax transformations.%
      \label{fig:comonoidal-adjunction-epsilon}%
    }
  \end{figure}
\end{example}

\section{Coends}\label{sec:coends}

\textsc{In this section we give the definition of} an enriched coend
in the special case of \(\Bbbk\)\hyp{}linear categories.
More general accounts can be found for example
in~\cite[Section~2.3]{gray80:closed},~\cite[Chapter~3]{kelly05:basic},
and~\cite[Chapter~4]{loregian2021}.

\begin{definition}\label{def:coend-vect}
  Let \(\cat{C}\) be an essentially small \(\Bbbk\)\hyp{}linear category.
  The \emph{coend} of a \(\Bbbk\)\hyp{}linear functor \(P\from \cat{C}^{\op} \otimes_{\Bbbk} \cat{C} \to \kVect\) is following the coequaliser in \(\kVect\):
  \[
    \begin{mytikzcd}[ampersand replacement=\&]
      \displaystyle{\coprod_{a, b \in \cat{C}} \cat{C}(a, b) \otimes_{\Bbbk} P(b, a)} \& \displaystyle{\coprod_{a \in \cat{C}} P(a,a)} \& \displaystyle{\int^{a \in \cat{C}} P(a,a).}
      \arrow[shift right=1, from=1-1, to=1-2]
      \arrow[shift left=1, from=1-1, to=1-2]
      \arrow[two heads, from=1-2, to=1-3]
    \end{mytikzcd}
  \]
  For \(a, b \in \cat{C}\) the two parallel morphisms are given by
  \begin{align*}
    \cat{C}(a,b)\otimes_{\Bbbk} P(b, a) &\to P(a,a), \qquad\qquad  f \otimes x \mapsto P(f, a)x, \\
    \cat{C}(a, b) \otimes_{\Bbbk} P(b, a) &\to P(b, b), \qquad\qquad g \otimes y \mapsto P(b,g)y.
  \end{align*}

  The \emph{end} \(\int_{a \in\cat{C}} P(a,a)\) of \(P\) is defined analogously as an equaliser in \(\kVect\).
\end{definition}

\begin{remark}
  A consequence of \(\cat{C}\) being essentially small is that we can rewrite the (co)equalisers defining the respective (co)ends in \cref{def:coend-vect} to index over a set of objects, instead of a proper class.
  Since \(\kVect\) is complete and cocomplete,
  the end and coend of any \(P\from \cat{C}^{\op} \otimes_{\Bbbk} \cat{C} \to \kVect\) exists.
\end{remark}

In case the indexing category \(\cat{C}\) can be inferred unambiguously from the context, we will omit writing it explicitly.

\begin{remark}
  Ends and coends have a number of useful properties,
  two of which we are going to frequently use in this thesis.

  First, they are functorial: Any natural transformation \(\alpha \from P \nt P'\) between functors
  \(P, P' \from\cat{C}^{\op} \times\cat{C} \to \kVect\) induces a unique morphism
  \[
    \int\!\mu \from \int\! P \to  \int\! P'
  \]
  by the universal property of (co-)equalisers.

  Second, for appropriate functors \(P\),
  the \emph{Fubini–Tonelli interchange law} holds:
  if any of the following coends exist:
  \[
    \int^{a \in \cat{C}} \int^{b \in \cat{C}}\! P(a, a, b, b), \quad
    \int^{(a, b) \in \cat{C} \times \cat{C}}\! P(a, a, b, b), \quad
    \int^{b \in \cat{C}} \int^{a \in \cat{C}}\! P(a, a, b, b),
  \]
  then they all do and are isomorphic;
  we often simply write \(\int^{a,b} P(a,a,b,b)\).
\end{remark}

The following lemma is also known as the \emph{enriched Yoneda lemma}.

\begin{lemma}\label{lem:enriched-yoneda}
  Let \(\cat{C}\) be a category,
  and let
  \(F \from \cat{C} \to \kVect\) and \(G \from \cat{C}^{\op} \to \kVect\) be functors.
  Then there are natural isomorphisms
  \begin{equation}\label{eq:coYoneda-covariant}
    \int^{c} \cat{C}(c, \blank) \otimes_{\Bbbk} Fc
    \;\cong\; F
    \;\cong\; \int_{c} \kVect(\cat{C}(\blank, c), Fc),
  \end{equation}
  and
  \begin{equation}\label{eq:coYoneda-contravariant}
    \int^{c} \cat{C}(\blank, c) \otimes_{\Bbbk} Gc
    \;\cong\; G
    \;\cong\; \int_{c} \kVect(\cat{C}(c, \blank), Gc).
  \end{equation}
\end{lemma}

As the name suggests,
the above lemma is an analogue of the Yoneda lemma in the enriched case;
as such, we will only refer to it as ``the Yoneda lemma'' in subsequent considerations.
For a proof, see~\cite[Section~2.4]{kelly05:basic}.

\begin{definition}\label{def:day-convolution}
  Let \(\cat{C}\) be a \(\Bbbk\)\hyp{}linear monoidal category.
  Given two \(\Bbbk\)\hyp{}linear functors \(F, G\from \cat{C}\to \kVect\),
  the \emph{Day convolution} of \(F\) and \(G\) is the \(\Bbbk\)\hyp{}linear functor
  \(F * G\from \cat{C} \to \kVect\) that is given by the coend
  \[
    F * G \defeq \int^{a, b} \cat{C}(a \otimes b, \blank) \otimes_{\Bbbk} Fa \otimes_{\Bbbk} Gb.
  \]
\end{definition}

\begin{remark}
  If \(\cat{C}\) is left or right closed monoidal,
  we may simplify the Day convolution product using the Yoneda lemma for coends:
  \begin{equation}\label{eq:day-convolution-simp-left}
    \begin{aligned}
      F * G &= \int^{a,b }\cat{C}(a \otimes b, \blank) \otimes_{\Bbbk} Fa \otimes_{\Bbbk} Gb
              \cong \int^{a,b }\cat{C}(a, {[b, \blank]}_{\ell}) \otimes_{\Bbbk} Fa \otimes_{\Bbbk} Gb \\
      \overset{\eqref{eq:coYoneda-covariant}}&{\cong} \int^{b } F({[b,\blank]}_{\ell}) \otimes_{\Bbbk} Gb,
    \end{aligned}
  \end{equation}
  \begin{equation}\label{eq:day-convolution-simp-right}
    \begin{aligned}
      F * G &= \int^{a,b }\cat{C}(a \otimes b, \blank) \otimes_{\Bbbk} Fa \otimes_{\Bbbk} Gb
              \cong \int^{a,b }\cat{C}(b, {[a, \blank]}_r) \otimes_{\Bbbk} Fa \otimes_{\Bbbk} Gb \\
      \overset{\eqref{eq:coYoneda-covariant}}&{\cong} \int^{a } Fa \otimes G({[a,\blank]}_r).
    \end{aligned}
  \end{equation}
\end{remark}

\begin{example}\label{ex:day-conv-as-tensor-product}
  Consider a \(\Bbbk\)\hyp{}linear monoidal category \(\cat{X}\) with a single object \(x\).
  Then \(A \defeq \mathrm{End}_{\cat{X}}(x)\) is a commutative algebra and \([\cat{X}, \kVect] \cong \rMod{A}\).
  Let \(F, G \from\cat{X} \to \kVect\) be two functors.
  Writing \(M \defeq Fx\) and \(N \defeq Gx\) for the corresponding modules over \(A\),
  and \(an \defeq na\) for all \(a\in A\) and \(n\in N\),
  we obtain the following using \cref{eq:day-convolution-simp-left} and the definition of coends:
  \[
    (F * G)x \cong \faktor{M \otimes_{\Bbbk} N}{\langle ma \otimes_{\Bbbk} n - m\otimes_{\Bbbk} an \mid m \in M, n \in N, a \in A\rangle} = M \otimes_{A} N.
  \]
  Thus, one recovers the tensor product of modules over commutative algebras.
\end{example}

\begin{theorem}[{\cite{day71:const}}]\label{thm:closed-mon-str-on-functors}
  For any monoidal category \(\cat{C}\),
  the category \([\cat{C}^{\op}, \kVect]\) is closed monoidal
  with Day convolution as its tensor product,
  \(\cat{C}(1, \blank)\) as its unit,
  and the internal homs given for all \(F, G\from \cat{C}^{\op} \to \kVect\) by
  \begin{align}
    {[F, G]}_{\ell} &\defeq \int_{a, b} \kVect(\cat{C}(\blank \otimes a, b), \kVect(Fa, Gb)),\label{eq:Day-internal-hom-left}\\
    {[F, G]}_r &\defeq \int_{a, b} \kVect(\cat{C}(a \otimes \blank, b), \kVect(Fa, Gb)).\label{eq:Day-internal-hom-right}
  \end{align}
\end{theorem}

\begin{remark}\label{rmk:day-closed-simplification}
  If the monoidal category \(\cat{C}\) is closed,
  then the formulas for the internal homs of \([\cat{C}, \kVect]\) may be simplified
  by means of the Yoneda lemma:
  \begin{equation}\label{eq:right-internal hom-simplification}
    \begin{aligned}
      &{[F,G]}_r
        = \int_{a,b} \kVect (\cat{C}(a \otimes \blank, b), \kVect(Fa, Gb)) \\
      &\quad \cong \int_{a, b} \kVect(\cat{C}(a \otimes \blank , b) \otimes_{\Bbbk}Fa, Gb)
        \cong \int_b \kVect\Big(\!\int^a\!\!\cat{C}(a \otimes \blank , b) \otimes_{\Bbbk}Fa , Gb\Big)\\
      &\quad \cong \int_b \kVect\Big(\!\int^a\!\!\!\cat{C}(a   ,{[\blank, b]}_{\ell}) \otimes_{\Bbbk}Fa , Gb\Big)
        \overset{\eqref{eq:coYoneda-covariant}}{\cong} \int_b \kVect(F{[\blank,b]}_{\ell}, Gb),
    \end{aligned}
  \end{equation}
  \begin{equation}\label{eq:left-internal hom-simplification}
    \begin{aligned}
      {[F,G]}_{\ell}
      \cong \int_b \kVect\Big(\int_a \cat{C}(a, {[\blank, b]}_r) \otimes_{\Bbbk}Fa , Gb\Big)
      \overset{\eqref{eq:coYoneda-covariant}}{\cong} \int_a \kVect(F{[\blank,b ]}_r, Gb).
    \end{aligned}
  \end{equation}
\end{remark}

\begin{remark}
  Whenever we treat \(\Cpsh \defeq [\cat{C}, \kVect]\) as a (closed) monoidal category,
  we implicitly equip with the convolution tensor product.
  Analogously, one may define a closed monoidal structure on \(\psh \defeq [\cat{C}^{\op}, \kVect]\).
  Note, however, that in this case we cannot simplify the internal hom in the same way as in \cref{rmk:day-closed-simplification};
  this would require the category \(\cat{C}^{\op}\), as opposed to \(\cat{C}\) itself, to be closed monoidal.
\end{remark}

The convolution structure is particularly well-behaved on representables:
\[
  \cat{C}(\blank, x) * \cat{C}(\blank, y)
  = \int^{a, b \in \cat{C}} \cat{C}(\blank, a \otimes b) \kotimes \cat{C}(a, x) \kotimes \cat{C}(b, y)
  = \cat{C}(\blank, x \otimes y)
\]
This connection extends to the entire functor category, see~\cite{day71:const,im86:univ-conv}.

\begin{proposition}\label{prop:yoneda-monoidal}
  For a monoidal category \(\cat{C}\),
  the Yoneda embedding
  \begin{equation}\label{eq:yoneda-embedding}
    \yo \from \cat{C} \to \psh = [\cat{C}^{\op}, \kVect], \qquad x \mapsto \cat{C}(\blank, x)
  \end{equation}
  is a strong monoidal functor.
\end{proposition}

\section{(Co)completions}

\textsc{In this subsection we give a brief—informal—account} of the results regarding the
monoidal pseudo\-functo\-riality of cocompletions
and the resulting (co)completion operations for monoidal and module categories.
We refer to~\cite{kelly05:basic,kashiwara06:categ} for generalities on (co)limits and (co)completions.
We implicitly assume all categories and functors to be \(\Bbbk\)\hyp{}linear.

Let \(\Phi\) be a class of diagrams.
We say that a category is \emph{\(\Phi\)\hyp{}cocomplete} if it admits colimits of functors with domain in \(\Phi\),
and we say that a functor is \emph{\(\Phi\)\hyp{}cocontinuous} if it preserves such colimits.

\begin{definition}\label{smoothcats}
  A monoidal category \(\cat{C}\) is called \emph{separately \(\Phi\)\hyp{}cocontinuous}
  if \(\cat{C}\) is \(\Phi\)\hyp{}cocomplete and its tensor product is separately \(\Phi\)\hyp{}cocontinuous.

  Similarly, for a \(\Phi\)\hyp{}cocomplete monoidal category \(\cat{D}\),
  a \(\cat{D}\)\hyp{}module category \(\cat{M}\) is said to be separately \(\Phi\)\hyp{}cocontinuous if \(\cat{M}\) is \(\Phi\)\hyp{}cocomplete
  and the action \(\blank \lact_{\cat{M}} \bblank\) is separately \(\Phi\)\hyp{}cocontinuous.
\end{definition}

Let \(\BiCat{Cat}_{\Cocts{\Phi}}\) be the \(2\)\hyp{}category of \(\Phi\)\hyp{}cocomplete categories,
\(\Phi\)\hyp{}cocontinuous functors,
and transformations between them.
Let \(\cat{C}\) be a small category.
By~\cite[Section~5.7]{kelly82:basic},
there is a category \(\Phi(\cat{C})\),
the \emph{\(\Phi\)\hyp{}cocompletion} of \(\cat{C}\),
and an embedding \(\iota\from \cat{C} \longhookrightarrow \Phi(\cat{C})\),
such that for any \(\Phi\)\hyp{}cocomplete category \(\cat{D}\)
the restriction functor \(\BiCat{Cat}_{\Cocts{\Phi}}(\cat{C},\cat{D}) \xrightarrow{\blank \circ \iota} \BiCat{Cat}(\cat{C},\cat{D})\) is an equivalence.

By the main results of~\cite{zoeberlein76:doctr,kock95:monad},
the inclusion \(2\)\hyp{}functor of \(\BiCat{Cat}_{\Cocts{\Phi}}\) into \(\BiCat{Cat}\) is \(2\)\hyp{}monadic,
with a left adjoint given by the pseudofunctor of \(\Phi\)\hyp{}cocompletion, \(\Phi(\blank)\from \BiCat{Cat} \to \BiCat{Cat}_{\Cocts{\Phi}}\).
This \(2\)\hyp{}monad is a lax-idempotent \(2\)\hyp{}monad in the sense of~\cite{kelly97}.

By~\cite{kelly00}, there is a similar \(2\)\hyp{}monad \(\Phi_{c}\) on \(\BiCat{Cat}\),
whose algebras are categories with a \emph{prescribed choice} of \(\Phi\)\hyp{}colimits.
Strict morphisms are functors strictly preserving the chosen \(\Phi\)\hyp{}colimits,
and pseudomorphisms are functors preserving \(\Phi\)\hyp{}colimits in the ordinary sense.

By~\cite[Theorem~6.2]{lopez11:pseud-kz}, this yields a pseudoclosed structure the \(2\)\hyp{}category \(\BiCat{Cat}_{\Cocts{\Phi_{c}}}\)
of categories with a choice of \(\Phi\)\hyp{}colimits, \(\Phi\)\hyp{}cocontinuous functors, and transformations between them.
The category underlying the internal hom from \(\cat{A}\) to \(\cat{B}\) is given by \(\BiCat{Cat}_{\Cocts{\Phi}}(\cat{A,B})\),
and the functors comprising the \(2\)\hyp{}monad are closed.
Further, in some cases, for instance if \(\Phi\) is the class of finite colimits,
the pseudoclosed structure becomes pseudoclosed monoidal.

Since monoidal categories and module categories can be formulated
as pseudomonoids internally to a pseudoclosed monoidal structure,
see~\cite{hyland02:pseud},
we find the following result.

\begin{proposition}\label{cocompletingnonsense}
  Let \(\cat{C}\) be a monoidal category.
  Then \(\Phi(\cat{C})\) is separately \(\Phi\)\hyp{}cocontinuous monoidal,
  the inclusion \(\iota\from \cat{C} \longhookrightarrow \Phi(\cat{C})\) is strong monoidal,
  and induces the following equivalence
  for any separately \(\Phi\)\hyp{}cocontinuous monoidal category \(\cat{D}\):
  \[
    \mathsf{StrMon}_{\Cocts{\Phi}}(\Phi(\cat{C}),\cat{D}) \xrightarrow{\blank \circ \iota} \mathsf{StrMon}(\cat{C,D}).
  \]
  Similar equivalences are induced for lax and oplax monoidal functors.

  Likewise, given a \(\cat{C}\)\hyp{}module category \(\cat{M}\),
  the \(\Phi\)\hyp{}cocompletion \(\Phi(\cat{M})\) is a separately \(\Phi\)\hyp{}cocontinuous \(\Phi(\cat{C})\)hyp{}module category,
  the inclusion \(\iota\from \cat{M} \longhookrightarrow \Phi(\cat{M})\) is a strong \(\cat{C}\)\hyp{}module functor, and the functor
  \begin{equation}\label{cocompletingmodules}
    \mathsf{Str}\Phi(\cat{C})\mathsf{Mod}_{\Cocts{\Phi}}(\Phi(\cat{M}),\cat{N})
    \xrightarrow{\blank \circ \iota}
    \mathsf{Str}\cat{C}\mathsf{Mod}(\cat{M}, \cat{N}),
  \end{equation}
  is an equivalence;
  similar equivalences exist for lax and oplax module functors.
\end{proposition}

We remark that explicit constructions of the monoidal and module structures obtained in \cref{cocompletingnonsense} and direct proofs of their universal properties
are given in a variety of cases and ways in the literature,
see for instance~\cite[Section~3]{mackaay19:simpl} and~\cite[Example~3.2.9]{capucci22:acteg-workin-amthem}.

\begin{proposition}
  Let \(\cat{A}\) be an additive category.
  The copower of \(\cat{A}\) over \(\kvect\)%
  —see \cref{ex:copower-as-comodule-monad}—%
  is the essentially unique \(\kvect\)\hyp{}module structure on \(\cat{A}\).
\end{proposition}

\begin{proof}
  Let \(\mathsf{add}\) denote the collection of finite discrete categories.
  An \(\mathsf{add}\)\hyp{}cocomplete category is simply an additive category, and an \(\mathsf{add}\)\hyp{}cocontinuous functor is an additive functor.

  Since any \(\Bbbk\)\hyp{}linear functor preserves direct sums, an additive \(\kvect\)\hyp{}module category is necessarily separately additive in the sense of Definition~\ref{smoothcats}.
  The result then follows by observing that \(\kvect \simeq \mathsf{add}(\{\oast\})\), where \(\{\oast\}\) is the terminal monoidal category, so
  \[
    \mathsf{StrMon}_{\mathsf{add}}(\kvect, \BiCat{Cat}_{\mathsf{add}}(\cat{A,A})) \simeq \mathsf{StrMon}(\{\ast\},\BiCat{Cat}_{\mathsf{add}}(\cat{A,A})) \simeq \{\ast\},
  \]
  where \(\{\ast\}\) is the terminal category and \(\mathsf{StrMon}_{\mathsf{add}}\) denotes the category of strong monoidal additive functors.
\end{proof}

\begin{proposition}\label{FinitaryAction}
  Let \(\cat{A}\) be an additive category admitting filtered colimits.
  The action of \(\kvect\) on \(\cat{A}\) extends uniquely
  to a separately finitary \(\kVect\)\hyp{}module category structure on \(\cat{A}\).
\end{proposition}
\begin{proof}
  Let \(\{\mathsf{add},\mathsf{filt}\}\) denote the collection of diagrams which are filtered or finite discrete.
  Then we have \(\kVect \simeq \{\mathsf{add},\mathsf{filt}\}(\{\oast\})\), and
  \begin{align*}
    &\mathbf{StrMon}_{\mathsf{filt}}(\kVect, \BiCat{Cat}_{\mathrm{add,filt}}(\cat{A,A}))
      = \mathbf{StrMon}_{\mathrm{add,filt}}(\kVect, \BiCat{Cat}_{\mathrm{add,filt}}(\cat{A,A}))\\
    &\simeq \mathbf{StrMon}(\{\oast\},\BiCat{Cat}_{\mathrm{add,filt}}(\cat{A,A}))
      = \mathbf{StrMon}(\{\oast\},\BiCat{Cat}_{\mathsf{filt}}(\cat{A,A})) \simeq \{\ast\}.\qedhere
  \end{align*}
\end{proof}

In view of \cref{lem:campus-bar},
filtered colimits and cocompletions under them,
which we shall call \emph{ind-completions},
are particularly important in the study of locally finite abelian categories.
For example, the ind-completion \(\Ind(\cat{C})\) of a locally finite abelian category \(\cat{C}\) always has enough injectives.
Another reason for our study of these completions
is that categories of comodules are locally finitely presentable,
whence we can make use of the characterisation of adjoint functors, see \cref{saftrex,saftlex} below.

\begin{definition}\label{def:lfp-category}
  An object \(x\) in a category \(\cat{C}\) is called \emph{compact}
  or \emph{finitely presented}
  if \(\cat{C}(x, \blank)\) is finitary; \ie, preserves filtered colimits.
  We denote the full subcategory of compact objects by \(\cat{C}_{\mathrm{c}}\).
  A category \(\cat{C}\) is called \emph{locally finitely presentable}
  if every object in \(\cat{C}\) is a filtered colimit of compact objects.
\end{definition}

Examples of ind-completions arise from the fact that
 \(\Ind(\tetramodfd[C][][][]) \simeq \Tetramod[C][][][]\).

\begin{proposition}[Special adjoint functor theorem for right adjoints]\label{saftrex}
  If \(\cat{C}\) and \(\cat{D}\) are locally finite abelian categories,
  then there are equivalences
  \[
    \mathsf{Rex}(\cat{C},\cat{D})
    \xrightarrow{\Ind} \mathsf{Cocont}(\Ind(\cat{C}), \Ind(\cat{D}))
    \iso \mathsf{Map}(\Ind(\cat{C}), \Ind(\cat{D})),
  \]
  where \(\mathsf{Map}(\cat{A,B})\) denotes the category of functors admitting a right adjoint.

  In particular, a functor from \(\Ind(\cat{C})\) to \(\Ind(\cat{D})\) admits a right adjoint if and only if it is cocontinuous,
  and the extension of a functor \(F\from \cat{C} \to \cat{D}\) to \(\Ind\)\hyp{}cocompletions is cocontinuous if and only if \(F\) is right exact.
\end{proposition}

\begin{proposition}[Special adjoint functor theorem for left adjoints]\label{saftlex}
  Let \(\cat{C}\) and \(\cat{D}\) be locally finite abelian categories.
  If a functor \(G\from \Ind(\cat{C}) \to \Ind(\cat{D})\) is continuous and preserves filtered colimits, then it admits a left adjoint.

  If a functor \(G\from \Ind(\cat{C}) \to \Ind(\cat{D})\) admits a left adjoint \(F\),
  then \(G\) preserves filtered colimits if and only if \(F\) preserves compact objects in \(\cat{C}\).
\end{proposition}

The latter part of \cref{saftrex} follows from \(\Ind(F)\) preserving both finite and filtered colimits, and thus preserving all colimits.
Since filtered colimits in a locally finitely presentable category are exact,
the extension \(\Ind(G)\) of a left exact functor \(G\) is left exact.
However, it is not clear that \(\Ind(G)\) preserves arbitrary products.

\begin{definition}
  An object \(X \in \Ind(\cat{C})\)
  for a locally finite abelian category \(\cat{C}\)
  is \emph{finitely cogenerated}
  if the functor \(\Ind(\cat{C})(\blank,X)\) preserves arbitrary products.
\end{definition}

One can verify that for a coalgebra \(D\) such that \(\Ind(\cat{C}) \simeq \rComod{D}\),
a comodule \(M\) is finitely cogenerated if and only if \(M\) embeds into a comodule of the form \(D^{\oplus m}\) for some \(m \in \nat\).
Observe that an object \(M \in \Ind(\cat{D})\) is compact if and only if it lies in the essential image of the embedding \(\cat{D} \longhookrightarrow \Ind(\cat{D})\).

\begin{definition}\label{qfFunctor}
  Let \(\cat{C}\) and \(\cat{D}\) be locally finite abelian categories.
  A functor \(F\from \Ind(\cat{C}) \to \Ind(\cat{D})\) is called \emph{quasi-finite} if
  for any finitely cogenerated injective object \(I\) in \(\Ind(\cat{C})\) and compact object \(M\) of \(\cat{D}\),
  the \(\Bbbk\)\hyp{}vector space \(\Ind(\cat{D})(M, FI)\) is finite\hyp{}dimensional.
\end{definition}

Using the equivalence \(\cat{C} \simeq \tetramodfd[C]{}\) of \cref{lem:campus-bar},
the following result is a direct consequence of~\cite[Proposition~1.3]{takeuchi77:morit}.

\begin{proposition}\label{locallyfiniteadjoints}
  Let \(G\from \cat{C} \to \cat{D}\) be a left exact functor of locally finite abelian categories.
  The functor \(\Ind(G)\) has a left adjoint if and only if it is quasi-finite.
\end{proposition}

A stronger result holds for finite abelian categories,
as a direct consequence of the Eilenberg--Watts theorem for categories of bimodules,
see for example~\cite[Lemma~2.1]{fuchs20:eilen-watts-radfor-s}.

\begin{proposition}\label{finiteadjoints}
  A functor \(F\from \cat{A} \to \cat{B}\) of finite abelian categories
  is right exact if and only if it has a right adjoint,
  and left exact if and only if it has a left adjoint.
\end{proposition}

Recall that an object \(A\) in a category \(\cat{A}\) that admits \(\Phi\)\hyp{}colimits is
said to be \emph{\(\Phi\)\hyp{}small} if the functor \(\cat{A}(A, \blank)\) preserves them.

\begin{lemma}\label{leftadjointssmallobjects}
  Let \(G\from \cat{C} \to \cat{D}\) be a \(\Phi\)\hyp{}cocontinuous functor,
  and assume \(G\) has a left adjoint \(F\).
  Then \(F\) sends \(\Phi\)\hyp{}small objects to \(\Phi\)\hyp{}small objects.
\end{lemma}
\begin{proof}
  Let \(x \in \cat{D}\) be a \(\Phi\)\hyp{}small object.
  Then
  \[
    \cat{C}(Fx,\blank) \cong \cat{D}(x,G(\blank)) = \cat{D}(x,\blank) \circ G.
  \]
  The functor \(\cat{C}(Fx,\blank)\) is thus naturally isomorphic to a composition of two \(\Phi\)\hyp{}cocontinuous functors,
  and hence it itself is a \(\Phi\)\hyp{}cocontinuous functor.
\end{proof}

%%% Local Variables:
%%% mode: LaTeX
%%% TeX-master: "../main"
%%% TeX-command-extra-options: "-shell-escape -fmt=prec.fmt -file-line-error -synctex=1"
%%% End:
